\chapter{The Theory of \texorpdfstring{$H^{p}$}{Hp} Spaces on Tubes}
\label{ch:c3}

Comparatively strong conditions have to be imposed on harmonic functions, defined in upper half-spaces, in order to insure the existence of boundary values. In sections two and three of the previous chapter we studied this situation. If, however, we consider a system of harmonic functions satisfying certain partial differential equations (for example, the real and imaginary parts of an analytic function form such a system satisfying the Cauchy--Riemann equations) weaker conditions are sufficient to insure the existence of boundary values. A more detailed discussion of these facts is presented in the first section. The rest of the chapter is devoted to the theory of holomorphic functions defined on tubes. These are, perhaps, the most natural domains where we can study the relation between the theory of holomorphic functions of several complex variables and Fourier analysis. In the second section we study the $H^2$ space theory of such functions. In the third section we consider an important special case of these spaces: the $H^2$ spaces over tubes whose bases are cones. An extension of the Paley--Wiener theorem to several variables is obtained in section four as a consequence of these results. In the fifth section we study the more general $H^p$ spaces over tube domains by reducing it to tubes based on certain special cones.

\section{Introductory Remarks}

We have seen in the previous chapter (see Theorem \ref{thm:c2-2-1}) how, by considering the Poisson integral of a function $f$ belonging to $L^p(E_n)$, $1 \leq p \leq \infty$, we obtain a harmonic function $u$, defined on $E_{n+1}^+$, satisfying

\begin{equation}\label{eq:c3-1-1}
\int\limits_{E_{n}}\left|u(x,y)\right|^{p} dx\leq A^{p}<\infty
\end{equation}

for all $y > 0$, where $A = \|f\|_p$. Furthermore, $u(x, y)$ converges non-tangentially to $f(x)$ for almost all $x \in E_n$ (see Theorem \ref{thm:c2-3-16}).

Theorem \ref{thm:c2-2-5} of the previous chapter tells us that inequalities of the type \eqref{eq:c3-1-1} characterize those harmonic functions on $E_{n+1}^{+}$ that are Poisson integrals, and, thus, imply the existence of nontangential boundary values. It can be shown that these results are best possible in the sense that if $0 < p < 1$ inequality \eqref{eq:c3-1-1} does not imply the existence of such boundary values as $y \to 0$; that is, for each such p there exists a harmonic function

satisfying \eqref{eq:c3-1-1} that does not have nontangential limits a.e. at the boundary, $E_{n}$, of $E_{n+1}^{+}$.

When n = 1 the situation changes very much if, instead of harmonic functions, we consider analytic functions. For example, if $F(z) = F(x + iy)$ is holomorphic for y > 0 (that is, in the upper half-plane $E_{2}^{+}$) and satisfies, for p > 0,

\begin{equation}\label{eq:c3-1-2}
\intop_{-\infty}^{\infty}|F(x+i y)|^{p}d x\leq A^{p}<\infty,\qquad y>0,
\end{equation}

then the limits $\lim_{z \to x} F(z) = F(x)$ exist at almost every $x \in (-\infty, \infty)$ if $z$ is allowed to approach $x$ nontangentially. The class of all holomorphic functions on $E_2^+$ satisfying \eqref{eq:c3-1-2} for some constant $A < \infty$ is denoted by $H^p = H^p(E_2^+)$. Better known are the corresponding $H^p$-spaces associated with the unit disk $D = \{z \in E_2; |z| < 1\}$. These consist of those analytic functions, $F$, in $D$ satisfying

\begin{equation}\label{eq:c3-1-3}
\intop_{-\pi}^{\pi}|F(r e^{i\theta})|^{p}d\theta\leq A^{p}<\infty,\qquad0\leq r<1.
\end{equation}

Again, the nontangential limits $F(e^{i\theta})$ exist for almost every $\theta \in [-\pi, \pi]$. Moreover, in both cases these are limits in the ``norm''; that is, when $F \in H^p(E_2^+)$ then

\begin{equation}\label{eq:c3-1-4}
\lim_{y\to0}\int\limits_{-\infty}^{\infty}|F(x+i y)-F(x)|^{p} dx=0,
\end{equation}

while, if $F \in H^{p}(D)$ then

\begin{equation}
\tag{1.4$'$}\label{eq:c3-1-4p}
\lim_{r\to1}\int\limits_{-\pi}^{\pi}|F(re^{i\theta})-F(e^{i\theta})|^{p}d\theta=0.
\end{equation}

It is our intention to extend these notions and results to higher dimensions. In so doing we meet with several essentially different approaches. One, involving the theory of functions of several complex variables, is the following. Let $B$ be an open subset of $E_n$. Then, the tube, $T_B$, with base\index{base (of tube)@base (of tube)} $B$ is the subset of all $z = (z_1, \ldots, z_n) = (x_1 + iy_1, \ldots, x_n + iy_n) = x + iy \in C_n$ (= $n$-dimensional complex Euclidean space) such that $y \in B$. For example, $E_2^+$ is the tube in $C_1$ with base\index{tube@tube} $B = \{y \in E_1; y > 0\}$. A function defined and holomorphic\footnote{We say that a function $F$ whose domain $D$ is contained in $C_n$ is holomorphic if for each $z_0 = (z_1^0, \ldots, z_n^0) \in D$ there exists a polydisk $\{z = (z_1, \ldots, z_n) \in C_n ; |z_1 - z_1^0| < r_1, \ldots, |z_n - z_n^0| < r_n\} \subset D$ in which $F$ is represented as an absolutely convergent power series: 
\[
F(z)=\sum_{k_{1},\cdots,k_{n}\geq0}a_{k_{1}\cdots k_{n}}(z_{1}-z_{1}^{0})^{k_{1}}\cdots(z_{n}-z_{n}^{0})^{k_{n}}.
\]} on a tube $T_B$ is said to belong to the space

$H^p = H^p(T_B), p > 0$, if there exists a constant $A < \infty$ such that

\begin{equation}\label{eq:c3-1-5}
\int\limits_{E_{n}}|F(x+i y)|^{p}d x\leq A^{p},\qquad\text{for all }y\in B.
\end{equation}

This clearly extends the definition of the spaces $H^p(E_2^+)$. The norm of $F \in H^p(T_B)$ is the positive number $\|F\|_p$ defined as the least $A$ for which \eqref{eq:c3-1-5} holds.

Another approach arises from the observation that a function $F = u + iv$ of one complex variable is analytic in a simply connected region if and only if $(v, u)$ is the gradient of a harmonic function in this region. Thus, we may consider a vector-valued function $F = (u_1, \ldots, u_n)$, defined in a region of $E_n$, to be a ``generalized analytic function'' if it is the gradient of a harmonic function in this region. More generally, we shall consider a vector-valued function $F = (u_1, u_2, \ldots, u_n)$, defined in a region $D \subset E_n$, for which the partial differential equations

\begin{equation}\label{eq:c3-1-6}
\sum_{j=1}^{n}\frac{\partial u_{j}}{\partial x_{j}}=0,\qquad\frac{\partial u_{i}}{\partial x_{j}}=\frac{\partial u_{j}}{\partial x_{i}},
\end{equation}

i, j = 1, 2,..., n, are satisfied. These conditions can also be stated in the form

\begin{equation}
\tag{1.6$'$}\label{eq:c3-1-6p}
\text{div}F=0,\qquad\text{curl}F=0.
\end{equation}

We note that the second condition implies that, in a simply connected subregion of $D$, $F$ is the gradient of some function $h$, while the first condition assures us that $h$ is harmonic. When $n = 2$ Eqs. \eqref{eq:c3-1-6} are, simply, the Cauchy--Riemann equations; in this case $u_2 + i u_1$ is an analytic function of $z = x_1 + ix_2$. If $F$ is such a function defined in $E_{n+1}^+$ and if there exists a constant $A < \infty$ such that

\begin{equation}\label{eq:c3-1-7}
\int\limits_{E_{n}}|F(x,y)|^{p} dx=\int\limits_{E_{n}}\left(\sum_{j=1}^{n+1}|u_{j}(x,y)|^{2}\right)^{p/2}dx\leq A^{p}
\end{equation}

for all $y > 0$,\footnote{We warn the reader that in \eqref{eq:c3-1-7} we are dealing with harmonic functions of $n + 1$ variables, where $y = x_{n+1}$.} we obtain an extension of the definition of the spaces $H^{p}(E_{2}^{+})$.

In this chapter we shall study the $H^{p}$ spaces associated with tube domains\index{Hp spaces for tube domains@$H^{p}$ spaces for tube domains}. The second extension will be studied in Chapter VI.

\section{The  $H^{2}$ Theory}

We fix an open connected subset, $B$, of $E_n$ and shall study the space $H^2(T_B)$. It is not hard to construct functions that belong to this space. In

fact, let f be a function satisfying

\begin{equation}\label{eq:c3-2-1}
\sup_{y\in B}\int\limits_{E_{n}}\left|f(t)\right|^{2}e^{-4\pi y\cdot t}dt\leq A^{2}<\infty.
\end{equation}

Let us write $z = x + iy$ and $z \cdot t = z_1t_1 + z_2t_2 + \cdots + z_n t_n$. We shall show that $|e^{2\pi i z \cdot t} f(t)| = e^{-2\pi y \cdot t} |f(t)|$ is majorized by an integrable function when y is restricted to a compact subset of $B$. We note that this implies

\begin{equation}\label{eq:c3-2-2}
F(z)=\int\limits_{E_{n}}e^{2\pi i z\cdot t}f(t)d t
\end{equation}

defines a holomorphic function in $T_{B}$.

It suffices to show that $e^{-2\pi y \cdot t}|f(t)|$ is majorized by an integrable function in a neighborhood of an arbitrary point $y_0 \in B$. Since $B$ is open there clearly exists such a neighborhood $N \subset B$; moreover,

\[
\int_{E_{n}}\bigl|f(t)\bigr|^{2}e^{-4\pi y_{0}\cdot t}e^{-4\pi(y-y_{0})\cdot t}\;d t\;\leq\;A^{2}
\]

whenever $y \in N$. Let us decompose $E_n$ into a finite union of nonoverlapping polygonal cones, $\Gamma_1, \Gamma_2, \ldots, \Gamma_k$, having vertex at the origin $0$ and such that whenever two points, $v$ and $w$, belong to one of these cones, then the angle between the segments $0v$ and $0w$ is less than, say, $\pi/4$. Since $N$ is a neighborhood of $y_0$ there exists $\delta > 0$ such that $\{y: |y - y_0| = \delta\} \subset N$. Let $\varepsilon = 4\pi\delta/\sqrt{2}$ and choose $y$ such that $(y_0 - y) \in \Gamma_j$ and $|y - y_0| = \delta$; then $\varepsilon|t| \leq -4\pi(y - y_0) \cdot t$ whenever $t \in \Gamma_j$. Consequently,

\[
\intop_{\Gamma_{f}}\big|f(t)\big|^{2}e^{-4\pi y_{0}\cdot t}e^{\varepsilon|t|}d t\;\leq\intop_{\Gamma_{f}}\big|f(t)\big|^{2}e^{-4\pi y_{0}\cdot t}e^{-4\pi(y-y_{0})\cdot t}d t\;\leq\;A^{2}.
\]

Hence,

\[
\int_{E_{n}}|f(t)|^{2}e^{-4\pi y_{0}\cdot t}e^{\varepsilon|t|}d t=\sum_{j=1}^{k}\int_{\Gamma_{j}}|f(t)|^{2}e^{-4\pi y_{0}\cdot t}e^{\varepsilon|t|}d t\leq k A^{2}<\infty.
\]

Therefore,

\[
\begin{aligned}&\int\limits_{E_{n}}|f(t)|e^{-2\pi y_{0}\cdot t}e^{(\varepsilon/4)|t|}dt\\=&\int\limits_{E_{n}}(|f(t)|e^{(\varepsilon/2)|t|}e^{-2\pi y_{0}\cdot t})e^{-(\varepsilon/4)|t|}dt\\\leq&\left(\int\limits_{E_{n}}|f(t)|^{2}e^{\varepsilon|t|}e^{-4\pi y_{0}\cdot t}dt\right)^{1/2}\left(\int\limits_{E_{n}}e^{-(\varepsilon/2)|t|}dt\right)^{1/2}<\infty.\end{aligned}
\]

Hence, for $y$ lying in a sphere of radius $\varepsilon/8\pi$ about $y_0$ we have

\[
|f(t)|e^{-2\pi y\cdot t}\le|f(t)|e^{-2\pi y_{0}\cdot t}e^{(\varepsilon/4)|t|},
\]

which is an integrable function of $t$.

An immediate application of the Plancherel theorem then gives us, for $y \in B$,

\[
\int_{E_{n}}|F(x+i y)|^{2}\,d x=\int_{E_{n}}|f(t)|^{2}e^{-4\pi y\cdot t}\,d t\leq A^{2}<\infty.
\]

Thus, the analytic function $F(z)$ defined by \eqref{eq:c3-2-2} belongs to $H^{2}(T_{B})$. The basic representation theorem concerning functions in this space asserts that all such functions have this form. That is,

\begin{theorem}
\label{thm:c3-2-3}
F belongs to $H^{2}(T_{B})$ if and only if it has the form \eqref{eq:c3-2-2} where f is a function satisfying \eqref{eq:c3-2-1}.
\end{theorem}

We shall illustrate the importance of this theorem by first deriving several of its consequences. We postpone its proof until the end of this section.

If $B \subset E_n$ we let $B^c$ denote the convex hull\index{convex hull@convex hull} of $B$; that is, $B^c$ is the smallest convex set containing $B$. It is easy to see that $B^c$ consists of all the finite sums $x = \sum \lambda_i x_i$, where $\sum \lambda_i = 1$, $\lambda_i \geq 0$ and $x_i$ belongs to $B$. It follows immediately from this that if $B$ is open so is $B^c$. The following corollary shows that we might as well restrict ourselves to bases $B$ that are convex.

\begin{corollary}
\label{cor:c3-2-4}
If $F \in H^2(T_B)$ then the integral \eqref{eq:c3-2-2} is well-defined for all $z \in T_{B^c}$ and gives us a function in $H^2(T_{B^c})$ having the same norm as $F$.
\end{corollary}

\textsc{Proof.}\ We first note that the Plancherel theorem, together with Theorem \ref{thm:c3-2-3}, implies that

\begin{equation}\label{eq:c3-2-5}
\|F\|_{2}=\sup_{y\in B}\left(\int_{E_{n}}|f(t)|^{2}e^{-4\pi y\cdot t}d t\right)^{1/2}.
\end{equation}

Let

\[
S=\Bigg\{y\in E_{n};\int\limits_{E_{n}}|f(t)|^{2}e^{-4\pi y\cdot t}d t\leq\|F\|_{2}^{2}\Bigg\}.
\]

Clearly, $B \subset S$ and it suffices to show that $S$ is convex. Suppose, therefore, that $y'$ and $y''$ belong to $S$ and $y = \alpha y' + (1 - \alpha)y''$, where $0 \leq \alpha \leq 1$. Using the inequality $u^{\alpha}v^{1 - \alpha} \leq \alpha u + (1 - \alpha)v$ (valid for any two non-

negative numbers u and v) we obtain

\[
\begin{aligned}\int\limits_{E_{n}}|f(t)|^{2}e^{-4\pi y\cdot t}dt&=\int\limits_{E_{n}}|f(t)|^{2}e^{-4\pi\alpha y^{\prime}\cdot t}e^{-4\pi(1-\alpha)y^{\prime \prime}\cdot t}dt\\&=\int\limits_{E_{n}}(|f(t)|^{2}e^{-4\pi y^{\prime}\cdot t})^\alpha(|f(t)|^{2}e^{-4\pi y^{\prime \prime}\cdot t})^{1-\alpha}dt\\&\leq\alpha\int\limits_{E_{n}}|f(t)|^{2}e^{-4\pi y^{\prime}\cdot t}dt+(1-\alpha)\int\limits_{E_{n}}|f(t)|^{2}e^{-4\pi y^{\prime \prime}\cdot t}dt\\&\leq\|F\|_{2}^{2}.\end{aligned}
\]

Thus $y \in S$, which shows that $S$ is convex. This proves the corollary.

From now on we shall assume the base $B$ to be convex, as well as open.

\begin{corollary}
\label{cor:c3-2-6}
A necessary and sufficient condition for $H^{2}(T_{B})$ to contain a function that is not identically zero is that no entire straight line lies in B.
\end{corollary}

\textsc{Proof.}\ Suppose $B$ contains the line consisting of all points $y$ satisfying $y = a\tau + b, -\infty < \tau < \infty$. Let $N(t_0)$ be a spherical neighborhood in $E_n$ about $t_0$ in which $a \cdot t$ is bounded away from 0. Then, for $y$ on this line, $F \in H^2(T_B)$ and $f$ satisfying \eqref{eq:c3-2-2} and \eqref{eq:c3-2-1} (the existence of such an $f$ is guaranteed by Theorem \ref{thm:c3-2-3})

\[
\|F\|_{2}^{2}\geq\int\limits_{E_{n}}|f(t)|^{2}e^{-4\pi y\cdot t}d t\geq\int\limits_{N(t_{0})}|f(t)|^{2}e^{-4\pi\tau(a\cdot t)}e^{-4\pi(b\cdot t)}d t.
\]

But, since we can choose $\tau$ arbitrarily on the real line, we can make the term $e^{-4\pi\tau(a\cdot t)}$ as large as we wish in $N(t_0)$. This shows that $f(t)=0$ for almost every $t$ in $N(t_0)$ and the first part of the corollary follows.

In order to see that if $B$ contains no line then $H^2(T_B)$ contains a function $F \neq 0$ we first observe that such a set $B$ implies the existence of an open convex cone, $\Gamma$, containing it and such that $\Gamma$ contains no entire straight lines.\footnote{We sketch a proof of this geometric fact in \eqref{frp:c3-6-15} below.} It follows that $\Gamma$ is regular (see §3), and for such cones, $H^2(T_T)$ does contain an $F \neq 0$. This, therefore, will establish Corollary \ref{cor:c3-2-6}.

In view of the one-dimensional theory we would expect a central question in the theory of $H^{p}$ spaces over a tube $T_{B}$ to be the following one: If $y_{0}$ is a boundary point of B does the limit

\begin{equation}\label{eq:c3-2-7}
F(x+iy_{0})=\lim\limits_{y\to y_{0},y\in B}F(x+iy)
\end{equation}

exist and, if so, in what sense?

A positive, but not altogether satisfactory, answer to this question can

be given. We first note that if $y_{0}$ is such a boundary point of B, Fatou's lemma and \eqref{eq:c3-2-5} imply

\[
\int_{E_{n}}|f(t)|^{2}e^{-4\pi y_{0}\cdot t}dt\leq\|F\|_{2}^{2}.
\]

Thus, since $f(t)e^{-2\pi y_0\cdot t}$ defines a function in $L^2(E_n)$, we can extend the definition of the integral \eqref{eq:c3-2-2} to $z = x + iy_0$ by taking the inverse Fourier transform of $f(t)e^{-2\pi y_0\cdot t}$. That is, written formally, we obtain the almost everywhere defined $L^2$-function of $x$

\begin{equation}\label{eq:c3-2-8}
F(x+i y_{0})=\int\limits_{E_{n}}e^{2\pi i(x+i y_{0})\cdot t}f(t)\,dt.
\end{equation}

In this sense, by taking the Fourier transforms of the functions $F_y = F(\cdot + iy)$, passing to the limit and then taking the inverse Fourier transform we can give meaning to the limit in \eqref{eq:c3-2-7}.

It is not unreasonable, however, to expect, as is the case in one dimension, that $F(x + iy)$ tends to $F(x + iy_0)$ as $y \in B$ tends to $y_0$, either in the $L^2$-norm or for almost all x. We shall now show, however, that this is not the case in general.

Suppose that $l$ is a line in $E_2$ and $y_1$ a point not on it. The equation of $l$ can be written in the form $y \cdot a = \beta$ for some fixed vector $a$ and real number $\beta$. $E_2$ is then split into two disjoint half-spaces: the set of all $y$ such that $y \cdot a > \beta$ and the set of all $y$ such that $y \cdot a < \beta$. Suppose $y_1$ belongs to the former. Then the function of two complex variables $z = (z_1, z_2) = (x + iy)$ defined by $G(z) = \exp \{-i\rho(z \cdot a - i\beta)\}$, $\rho > 0$, satisfies $|G(z)| = \exp \{\rho(y \cdot a - \beta)\}$, which is 1 for $y$ on $l$, less than one in the half-space not containing $y_1$ and is equal to a number $N > 1$ if $y = y_1$. This $N$ can be made as large as we wish by choosing $\rho$ large enough. Moreover, $|G(z)| \leq N$ if $z = x + iy$ satisfies $(y_1 - y) \cdot a \geq 0$.

Now, suppose $B$ is a disk in $E_2$, having 0 on its boundary and lying in the upper half-plane. Choose a sequence $\{y_k\}$ on the boundary of $B$ that converges to 0 and let $\{\sigma_k\}$ be an accompanying sequence of sectors of $B$, each of which consists of a region between a line $l_k$ cutting the boundary of $B$ on both sides of $y_k$ and the arc on this boundary that contains $y_k$ and is determined by these two points of intersection. We further assume that $l_k$ is so close to $y_k$ that the $\sigma_k$'s are mutually disjoint (see figure) and that $l_k$ is parallel to the tangent to $B$ at $y_k$. For each $k$ we can construct a function, $G_k$, of the type described in the previous paragraph, satisfying

(i) $G_{k}$ is analytic in $C_{2}$;

(ii) $|G_{k}(x + iy)|$ depends only on y;

(iii) $|G_k(z)| \leq 1$ if $z \in T_B - T_{\sigma_k}$;

\[
{\left|G_{k}(x+i y_{k})\right|}=1+2^{k+2}=N_{k}\text{ and, for }z\in T_{\sigma_{k}},\left|G_{k}(z)\right|\leq N_{k}.
\]

\begin{figure}[htbp]
\centering
\figtube
\caption{The tube $B$ and the boundary arcs $y_k$, $y_{k+1}$}
\label{fig:c3-p108}
\end{figure}

We then define a function $F$ by letting

\[
\begin{array}{c} F(z) = \displaystyle\sum_{k=1}^{\infty} 2^{-k}G_{k}(z).\end{array}
\]

If $z \in T_B$ then it either belongs to precisely one of the tubes $T_{\sigma_k}$, say $z \in T_{\sigma_{k_0}}$, or it belongs to none of them; in the first case, by (iii) and in the second part of (iv), we have

\[
\left|F(z)\right|\leq\sum_{k=1}^{k_{0}-1}2^{-k}+(2^{-k_{0}}+4)+\sum_{k=k_{0}+1}^{\infty}2^{-k}=5,
\]

while in the second case $|F(z)| \leq \sum_{k=1}^{\infty} 2^{-k} = 1$. In particular, this shows that $F \in H^{\infty}(T_B)$. By (ii) and (iv) we can find points $y_k' \in \sigma_k$ so close to $y_k$ that $|G_k(x + iy'_k)| > N_k - 1$, for $k = 1, 2, \ldots$. Thus,

\[
\begin{aligned}{|F(x+i y_{k}^{\prime})|\geq2^{-k}|G_{k}(x+i y_{k}^{\prime})|}&{{}-\sum_{j\neq k}2^{-j}|G_{k}(x+i y_{k}^{\prime})|}\\ {}&{{}>2^{-k}(N_{k}-1)-\sum_{j\neq k}2^{-j}>4-1=3.}\\ \end{aligned}
\]

This shows that if we approach the boundary point 0 with a sequence that avoids all sectors $\sigma_k$ we have $|F(x + iy)| \leq 1$ for all $y$ in this sequence. On the other hand $\lim_{k \to \infty} y'_k = 0$ and $|F(x + iy'_k)| > 3$. In particular, this shows that the pointwise limit

\[
\lim\limits_{y\in B,y\to0}F(x+iy)
\]

cannot exist for any $x \in E_{2}$.

To give an example of a function in $H^2(T_B)$ that does not have a limit in the $L^2$-norm as we approach 0 arbitrarily within $B$ it suffices to find a function $G\in H^2(T_B')$, where $\bar{B}\subset B$, such that $G(x+i0)=G(x)\not\equiv0$ and multiply it by the $F$ just obtained. Then $\int_{E_n}|F(x+iy)G(x+iy)|^2\,dx\geq9\int_{E_n}|G(x+iy)|^2\,dx$ whenever $y=y_k^!, k=1,2,\ldots$, while

\[
\int\limits_{E_{n}}|F(x+i y)G(x+i y)|^{2}d x\leq\int\limits_{E_{n}}|G(x+i y)|^{2}d x
\]

if $y \in B$ belongs to none of the $\sigma_k$. Thus, the limit

\[
\lim\limits_{y\to0}\boldsymbol{F}(x+iy)\boldsymbol{G}(x+iy)
\]

cannot exist in the $L^{2}$-norm. $G(z) = 1/(z_{1} + i)(z_{2} + i)$ for $z = (z_{1}, z_{2}) = (x_{1} + iy_{1}, x_{2} + iy_{2})$, $y_{1} > -\frac{1}{2}$, $y_{2} > -\frac{1}{2}$, is an example of such a function. It is clearly defined and analytic in a tube $T_{B^{\prime}}$ with $\overrightarrow{B} \subset B^{\prime}$; furthermore,

\[
\begin{aligned}\int\limits_{E_{2}}|G(x+iy)|^{2}dx&=\int\limits_{-\infty}^{\infty}\int\limits_{-\infty}^{\infty}\frac{1}{|x_{1}+i(y_{1}+1)|^{2}|x_{2}+i(y_{2}+1)|^{2}}dx_{1}dx_{2}\\&\leq\left\{\int\limits_{-\infty}^{\infty}\frac{1}{x^{2}+\frac{1}{4}}dx\right\}^{2}<\infty.\end{aligned}
\]

Thus, $G \in H^2(T_B')$

If we specialize $B$, however, we do obtain $L^{2}$-limits as we approach a boundary point. We define an open polyhedron\index{polyhedron (convex, open)@polyhedron (convex, open)} in $E_n$ to be the interior of the convex hull of a finite subset of $E_n$. The following result asserts, in particular, that when $B$ is an open polyhedron the limits \eqref{eq:c3-2-7} exist in the $L^{2}$-sense at all boundary points of $B$:

\begin{corollary}
\label{cor:c3-2-9}
Let $P$ be an open polyhedron in $E_n$ and $F \in H^2(T_P)$. If we then extend the definition of $F$ to the set $T_P$ by \eqref{eq:c3-2-8} the mapping $y \to F(x + iy)$, from $\overline{P}$ to $L^2(E_n)$, is continuous.
\end{corollary}

\textsc{Proof.}\ By the Plancherel theorem it is sufficient to show that the mapping $y \to f(t)e^{-2\pi y\cdot t}$ is continuous. Suppose $\bar{P}$ is the convex hull of

the finite set $\{y_1, y_2, \ldots, y_k\} \subset E_n$. Let

\[
G(t)=\sum_{j=1}^{k}e^{-4\pi y_{j}\cdot t}|f(t)|^{2}.
\]

The function $G$ is clearly integrable in $E_n$. Moreover, it majorizes $e^{-4\pi y \cdot t}|f(t)|^2$ for each $y$ in $\overline{P}$. For, if $y \in \overline{P}$, then $y = \alpha_1 y_1 + \alpha_2 y_2 + \cdots + \alpha_k y_k$, where the $\alpha_j$'s are nonnegative numbers whose sum is 1; but this implies that

\[
e^{-4\pi y\cdot t}=\exp\left\{-4\pi\sum_{j=1}^{k}\alpha_{j}(y_{j}\cdot t)\right\}=\prod_{j=1}^{k}(e^{-4\pi(y_{j}\cdot t)})^{\alpha_{j}}\;\leq\;\sum_{j=1}^{k}\alpha_{j}e^{-4\pi y_{j}\cdot t}.
\]

Now, for $y, \bar{y} \in \bar{P}, y \to \bar{y}$ we have

\[
|f(t)e^{-2\pi y\cdot t}-f(t)e^{-2\pi\bar{y}\cdot t}|^{2}\to0;
\]

moreover, this convergence is dominated by $4G(t)$. The corollary now follows from the Lebesgue dominated convergence theorem.

\begin{corollary}
\label{cor:c3-2-10}
Let $B$ be an open convex subset of $E_n$ and $y_0$ a point on its boundary. Suppose $F \in H^2(T_B)$ and $P$ is an open polyhedron contained in $B$ having $y_0$ as a boundary point, then $F(x + iy)$ tends to $F(x + iy_0)$ in the $L^2$-norm as $y$ tends to $y_0$ within $P$, where $F(x + iy_0)$ is the function defined by \eqref{eq:c3-2-8}.
\end{corollary}

Since $H^{2}(T_{B}) \subset H^{2}(T_{P})$ this is just a special case of Corollary \ref{cor:c3-2-9}.

It is not hard to see that the counter example we have just given can be modified in such a way as to include each boundary point that is not a polygonal boundary point\index{polygonal boundary point@polygonal boundary point}; by this term we mean a point $y_0$ that is the point of intersection of two straight line segments making up part of the boundary $\partial B$ of $B$ (we do not exclude the possibility that these two segments lie on the same line; in this case, $y_0$ would be interior to a single line segment that lies in $\partial B$). If $y_0$ is not a polygonal boundary point, it follows from the convexity of $B$ that we can find a sequence $\{y_i\}$ in $\partial B$ converging to $y_0$ and such that the open segments $\overline{y_i y_0}$ lie in the interior of $B$. This clearly allows us to find a sequence $\{\sigma_k\}$ of mutually disjoint sectors with accompanying analytic functions $\{G_k\}$ enabling us to construct a function $F$ just as we did in our counter example. This, together with Corollary \ref{cor:c3-2-9}, gives us the following necessary and sufficient conditions for the limits \eqref{eq:c3-2-7} to exist in the $L^2$ sense, when the dimension is two:

\begin{theorem}
\label{thm:c3-2-11}
If $B$ is an open convex subset of $E_2$ and $y_0 \in \partial B$ then
\[
\lim_{y\to y_{0},y\in B}F(x+iy)=F(x+iy_{0})
\]

exists in $L^{2}$ for every $F\in H^{2}(T_{B})$ if and only if $y_{0}$ is a polygonal boundary point of $B$.
\end{theorem}

In general, if the limits in \eqref{eq:c3-2-7} exist (in the $L^{2}$ sense, or for almost every $x$, or in $L^{p},\ldots$) we shall refer to them as unrestricted limits\index{restricted limits@restricted limits!in tubes@in tubes} (in $L^{2}$, or a.e., or in $L^{p},\ldots$). If such limits exist whenever $y$ converges to $y_{0}$ within a polyhedron in $B$ having $y_{0}$ as a boundary point we say that $(L^{2}, \text{ or a.e.},\text{ or }L^{p},\ldots)$ restricted limits exist at $y_{0}$. In this terminology, Corollary \ref{cor:c3-2-10} asserts that $L^{2}$ restricted limits exist at all boundary points for all $H^{2}$ functions; Theorem \ref{thm:c3-2-11}, on the other hand, gives necessary and sufficient conditions, when $n = 2$, for the existence of $L^{2}$ unrestricted limits\index{unrestricted limits@unrestricted limits}.

We now pass to the proof of Theorem \ref{thm:c3-2-3}. We begin with a lemma asserting, in particular, that if $F \in H^2(T_B)$ then $F$ is uniformly bounded in ``strictly smaller'' tubes. We state and prove this lemma for the wider class of spaces $H^p(T_B)$ (see \eqref{eq:c3-1-5}) since we will need this more general result later on. More precisely,

\begin{lemma}
\label{lem:c3-2-12}
Suppose $F \in H^p(T_B)$, $p > 0$, and $B_0 \subset B$ satisfies $d(B_0, B^{\circ}) = \inf\{|y_1 - y_2|; y_1 \in B_0, y_2 \notin B\} \geq \varepsilon > 0$ then there exists a constant $C = C(\varepsilon, n)$, depending on $\varepsilon$ and $n$ but not on $F$, such that
\[
\sup_{z\in T_{B_{0}}}|F(z)|\;\leq\;C\|F\|_{p}.
\]

\end{lemma}

\textsc{Proof.}\ Let $z_0 = x_0 + iy_0 \in T_{B_0}$ and $S_\varepsilon = \{z \in C_n; |z - z_0| < \varepsilon\}$. If $\Sigma_\varepsilon = \{y \in E_n; |y - y_0| < \varepsilon\}$ then $S_\varepsilon \subset T_{\Sigma_\varepsilon} \subset T_B$. Hence, recalling that $\Omega_m$ denotes the volume of the unit ball of $E_m$, $m \geq 1$,

\[
\begin{aligned}{\left(\int_{S_{\varepsilon}}|F(z)|^{p}d x d y\right)^{1/p}}&{{}\leq\left(\int_{T_{\Sigma_{\varepsilon}}}|F(z)|^{p}d x d y\right)^{1/p}}\\ {}&{{}=\left(\int_{\Sigma_{\varepsilon}}\left\{\int_{E_{n}}|F(x+i y)|^{p}d x\right\}d y\right)^{1/p}}\\ {}&{{}\leq\|F\|_{p}(\Omega_{n}\varepsilon^{n})^{1/p}.}\\ \end{aligned}
\]

On the other hand, since $|F|^p$ is subharmonic as a function of 2n variables (see example (3) in Section 4 and Theorem \ref{thm:c2-4-4p} in Chapter II),

\[
|F(z_{0})|^{p}\leq\Omega_{2n}^{-1}\varepsilon^{-2n}\Biggl(\int_{S_{\varepsilon}}|F(z)|^{p}d x d y\Biggr).
\]

Combining these inequalities we obtain $|F(z_0)| \leq C\|F\|_p$, with $C = (\Omega_n/\Omega_{2n})^{1/p} \varepsilon^{-n/p}$.

In order to establish Theorem \ref{thm:c3-2-3} we need to show that if $F \in H^2(T_B)$ then there exists a function $f$ satisfying \eqref{eq:c3-2-1} giving us the representation \eqref{eq:c3-2-2} for $F$ (the converse has already been established).

Let $f_y$, for $y \in B$, be the Fourier transform of $F(x + iy)$ considered as a function of $x$. It clearly suffices to show that if $y, y' \in B$ then $e^{2\pi y' \cdot t} f_y(t) = e^{2\pi y \cdot t} f_y(t)$, a.e. in $t$. For, then, $f(t) = e^{2\pi y \cdot t} f_y(t)$ is almost everywhere

defined independently of $y \in B$, $f$ satisfies \eqref{eq:c3-2-1} with $A = \| F\|$ and \eqref{eq:c3-2-2} must hold (see the argument preceding Theorem \ref{thm:c3-2-3}). It clearly suffices to show this when $y$, $y'$ belong to a cube $Q$ whose closure lies in $B$ and whose sides are parallel to the coordinate axes. We further assume, for the moment, that $|F(x + iy)|$, as a function of $x$, is majorized by a function vanishing sufficiently rapidly at $\infty$, uniformly in $y \in Q$. Lastly, we remark that, by an iteration argument, it suffices to consider the case where the two points $y$, $y' \in Q$ have the form $y = (\eta_1, y_2, \ldots, y_n)$ and $y' = (\eta_1', y_2, \ldots, y_n)$. Suppose $\eta_1' \geq \eta_1$, then by the Cauchy integral theorem,

\[
\begin{aligned}0=&\int\limits_{-R}^{R}e^{-2\pi i(x_{1}+\dot{i}\eta_{1})t_{1}}F(x_{1}+\dot{i}\eta_{1},\ldots,x_{n}+\dot{i}y_{n})dx_{1}\\&+\int\limits_{\eta_{1}}^{\eta_{1}^{\prime}}e^{-2\pi i(R+\dot{i}\eta)t_{1}}F(R+\dot{i}\eta,\ldots)d\eta\\&+\int\limits_{R}^{-R}e^{-2\pi i(x_{1}+\dot{i}\eta_{1}^{\prime})t_{1}}F(x_{1}+\dot{i}\eta_{1}^{\prime},\ldots)dx_{1}\\&+\int\limits_{\eta_{1}^{\prime}}^{\eta_{1}}e^{-2\pi i(-\dot{R}+\dot{i}\eta)t_{1}}F(-R+\dot{i}\eta,\ldots)d\eta.\end{aligned}
\]

Because of our assumptions, the 2nd and 4th terms tend to 0 as $R \to \infty$. Thus, after an integration in $x_2, \ldots, x_n$, we get

\[
\begin{aligned}{e^{2\pi y\cdot t}f_{y}(t)}&{{}=\intop_{E_{n}}e^{-2\pi i(x+i y)\cdot t}F(x+i y)\;d x}\\ {}&{{}=\intop_{E_{n}}e^{-2\pi i(x+i y^{\prime})\cdot t}F(x+i y^{\prime})\;d x}\\ {}&{{}=e^{2\pi y^{\prime}\cdot t}f_{y^{\prime}}(t).}\\ \end{aligned}
\]

We now show how the case of the general $F\in H^2(T_B)$ follows from this one. By Lemma \ref{lem:c3-2-12} $F$ is bounded on $T_Q$, say $|F(z)|\leq M$ for $z\in T_Q$. Thus, if we define $F^{(\varepsilon)}$, for $\varepsilon>0$, by letting

\[
\begin{array}{r}{F^{(\varepsilon)}(z)=e^{\left\{-\varepsilon\displaystyle\sum_{j=1}^{n}z_{j}^{2}\right\}}F(z),}\end{array}
\]

for $z = (z_1, \ldots, z_n) \in T_Q$, we have

\[
|F^{(\varepsilon)}(x+i y)|\leq M e^{\varepsilon n a^{2}}e^{-|x|^{2}\varepsilon}
\]

for $y \in Q$, where $a = \max_{y \in Q} \{|y_1|, \ldots, |y_n|\}$. The argument just presented,

therefore, applies to the Fourier transforms $f_{y}^{(\varepsilon)}$ of $F^{(\varepsilon)}(x + iy)$ and we have

\begin{equation}\label{eq:c3-2-13}
e^{2\pi y\cdot t}f_{y}^{(\varepsilon)}(t)=e^{2\pi y^{\prime}\cdot t}f_{y^{\prime}}^{(\varepsilon)}(t)
\end{equation}

for $y, y' \in Q$. But $\int_{E_n} |F^{(\varepsilon)}(x + iy) - F(x + iy)|^2 \, dx \to 0$ as $\varepsilon \to 0$. Thus, the equality \eqref{eq:c3-2-13} is preserved in the limit if we let $\varepsilon \to 0$. That is, we can apply the Plancherel theorem to obtain the $L^2$-convergences $f_y^{(\varepsilon)} \to f_y$ and $f_y^{(\varepsilon)} \to f_{y'}$. In particular

\[
e^{2\pi y^{\prime}\cdot t}f_{y}(t)=e^{2\pi y\cdot t}f_{y^{\prime}}(t)
\]

for almost every t. This proves the theorem.

\section{Tubes Over Cones}

It is natural to expect that, if we restrict the type of base, the theory of $H^2$ spaces (or, more generally, $H^p$ spaces) becomes richer. This is particularly true if we consider tubes over bases that are open cones. By the latter we mean nonempty subsets $\Gamma \subset E_n$ satisfying (i) $0 \notin \Gamma$, and (ii) whenever $x, y \in \Gamma$ and $\alpha, \beta > 0$ then $\alpha x + \beta y \in \Gamma$. We note that, in particular, $\Gamma$ is convex.

A closed cone is the closure of an open cone. It is clear that if $\Gamma$ is an open cone then $\Gamma^* = \{x \in E_n; x \cdot t \geq 0, t \in \Gamma\}$ is closed. Moreover, if $\Gamma^*$ has a nonvoid interior it is a closed cone; in this case we say that $\Gamma$ is a regular cone. $\Gamma^*$ is called the cone dual to $\Gamma$. When $n = 1$ the open cones are the half-lines $\{x \in E_1; x > 0\}$ and $\{x \in E_1; x < 0\}$. When $n = 2$ the open cones are the angular regions between two rays meeting at the origin and forming an angle that is less than or equal to $\pi$. Such a cone will be regular if and only if this angle is strictly less than $\pi$.\index{cone@cone!dual@dual}\index{cone@cone!regular@regular}

When the base is a cone one obtains the following representation theorem for $H^{2}$-functions (a sharper result than Theorem \ref{thm:c3-2-3}):

\begin{theorem}
\label{thm:c3-3-1}
Suppose $\Gamma$ is an open cone. Then $F \in H^2(T_\Gamma)$ if and only if
\begin{equation}\label{eq:c3-3-2}
F(z)=\int\limits_{\Gamma^{\bullet}}e^{2\pi i z\cdot t}f(t)d t
\end{equation}

where f is a measurable function on $E_{n}$ satisfying
\[
\int\limits_{\Gamma^{\bullet}}|f(t)|^{2}d t<\infty.
\]

\end{theorem}

Furthermore,

\[
\|F\|_{2}=\left(\int_{\Gamma^{*}}|f(t)|^{2}d t\right)^{1/2}.
\]

Consequently, the correspondence $F \leftrightarrow f$ is a unitary linear mapping of

$H^2(T_\Gamma)$ onto $L^2(\Gamma^*)$. In particular, $H^2(T_\Gamma)$ contains a function that is not identically zero if and only if $\Gamma$ is regular.

\textsc{Proof.}\ Since $y \cdot t \geq 0$ whenever $y \in \Gamma$ and $t \in \Gamma^*$, the fact that $F \in H^2(T_\Gamma)$ when it has the representation \eqref{eq:c3-3-2}, with $f \in L^2(\Gamma^*)$, is a consequence of Theorem \eqref{thm:c3-2-3}.

Conversely, suppose $F \in H^2(T_\Gamma)$. Then, by Theorem \ref{thm:c3-2-3} and \eqref{eq:c3-2-5},

\[
F(z)=\int\limits_{E_{n}}e^{2\pi i z\cdot t}f(t)d t,
\]

where

\[
\left\|F\right\|_{2}^{2}=\sup_{y\in\Gamma}\int\limits_{E_{n}}e^{-4\pi y\cdot t}|f(t)|^{2}dt.
\]

We shall show that $f$ vanishes a.e. in the complement of $\Gamma^*$. If $t_0 \notin \Gamma^*$ then there exists $y_0 \in \Gamma$ such that $y_0 \cdot t_0 < 0$. Thus, there exists a neighborhood of $t_0$, $N = N(t_0) \subset E_n - \Gamma^*$, and a $\delta > 0$ such that $y_0 \cdot t < -\delta < 0$ for $t \in N$. Hence, $(ky_0) \cdot t < -k\delta$ for $t \in N$ and positive $k$. Since $ky_0 \in \Gamma$ we also have $\int_N e^{-4\pi k y_0 \cdot t} |f(t)|^2 \, dt \leq \int_{E_n} e^{-4\pi k y_0 \cdot t} |f(t)|^2 \, dt \leq \|F\|_2^2 < \infty$. But this implies that

\[
\int\limits_{N} e^{4\pi k\delta}|f(t)|^{2}\,dt\ \leq\ \|F\|_{2}^{2}<\infty
\]

for all $k > 0$, and this clearly implies that $f(t) = 0$ for almost every $t$ in $N(t_0)$. It follows that $f'(t) = 0$ for almost all $t$ outside $\Gamma^*$.\footnote{Observe that this part of the proof is very similar to that of Corollary \ref{cor:c3-2-6}.} The other parts of the statement of Theorem \ref{thm:c3-3-1} are now obvious.

In the classical case, when $F$ belongs to the $H^{2}$ space associated with the upper half plane (i.e., when the base is the cone of positive real numbers) it is easy to establish the Cauchy integral representation

\begin{equation}\label{eq:c3-3-3}
F(z)=\frac{1}{2\pi i}\int\limits_{E_{1}}\frac{F(\zeta)}{\zeta-z}d\zeta
\end{equation}

of $F$ in terms of its boundary values. We shall now show that this representation can be extended to include the case when the base is a cone, $\Gamma$, in $E_n$; moreover, this representation involves only those boundary values obtained as $y \in \Gamma$ tends to the vertex of $\Gamma$ (i.e., the origin, $0$, of $E_n$). The fact that these boundary values exist in the $L^2$ sense is an easy consequence of Theorem \ref{thm:c3-3-1}:

\begin{corollary}
\label{cor:c3-3-4}
Suppose $\Gamma$ is an open cone in $E_n$ and $F(x + iy) \in H^2(T_\Gamma)$ then there exists a function $F(x)$, defined on $E_n$, such that $F(x + iy) \to F(x)$ in the $L^2$-norm as $y \in \Gamma$ tends to 0.
\end{corollary}

\textsc{Proof.}\ Using the notation of Theorem \ref{thm:c3-3-1} we define $F(x)$ to be the inverse Fourier transform of $f$ (which we may think of as defined on $E_n$ and 0 outside of $\Gamma^*$); thus, formally,

\[
F(x)=\int\limits_{\Gamma^{*}}e^{2\pi i x\cdot t}f(t)\;d t.
\]

But \eqref{eq:c3-3-2} asserts that $F(x + iy)$ is the inverse Fourier transform of $e^{-2\pi y\cdot t}f(t)$; therefore, since the latter converges to $f(t)$ in the $L^2$-norm as $y \to 0$, the desired conclusion follows from the Plancherel theorem.

The extension of \eqref{eq:c3-3-3} described above will be obtained by means of the following kernel function associated with a tube domain: For $z = x + iy$ in $T_{r}$ we let

\[
K(z)=\int\limits_{\Gamma^{*}}e^{2\pi i z\cdot t}\;d t=\int\limits_{\Gamma^{*}}e^{2\pi i x\cdot t}e^{-2\pi y\cdot t}\;d t.\footnote{When $\Gamma$ is the cone of positive real numbers then $K(z) = -1/2\pi iz$ and \eqref{eq:c3-3-3} has the form $F(z) = \int_{E_1} K(z - \xi) F(\xi) \, d\xi$. Theorem \ref{thm:c3-3-6} shows that this formula holds for functions in a general tube $T_\Gamma$ whose base is a cone.}
\]

We thus obtain a continuous function $K$ defined on $T_{\Gamma}$, called the Cauchy kernel\index{Cauchy kernel@Cauchy kernel} associated with the tube $T_{\Gamma}$. As a function of $x = \text{Re}\{z\}$, $K$ belongs to $L^{2}(E_{n})$; in fact, it follows from Plancherel's theorem that

\begin{equation}\label{eq:c3-3-5}
\int_{E_{n}}|K(x+i y)|^{2}\,d x=\int_{\Gamma^{\bullet}}e^{-4\pi y\cdot t}\,d t=K(2i y)
\end{equation}

for all $y \in \Gamma$. (We shall see momentarily that $K(2iy)$ is finite.)

\begin{theorem}
\label{thm:c3-3-6}
If $F \in H^2(T_\Gamma)$ then
\[
F(z)=\int\limits_{E_{n}}K(z-\xi)F(\xi)d\xi
\]

for all $z \in T_{\Gamma}$, where $F(\xi) = \lim_{\eta \to 0, \eta \in \Gamma} F(\xi + i\eta)$ is the limit function of Corollary \ref{cor:c3-3-4}.
\end{theorem}

\textsc{Proof.}\ From Theorem \ref{thm:c3-3-1} and 3.4 we have

\[
F(z)=F(x+i y)=\int\limits_{\Gamma^{\bullet}}e^{2\pi i z\cdot t}f(t)d t
\]

where $f$ is the Fourier transform of $F(\xi) = \lim_{\eta \to 0, \eta \in \Gamma} F(\xi + i\eta)$. That is, $f$ is the limit in the $L^2$ norm of the sequence of functions

\[
f_{k}(t)\;=\;\int\limits_{\{\xi\}\leq k}F(\xi)e^{-2\pi i t\cdot\xi}\;d\xi,
\]

$k = 1, 2, 3, \ldots$ Consequently, using Fubini's theorem and the fact that $K(x - \xi + iy)$, as a function of $\xi$, belongs to $L^2(E_n)$ (see \eqref{eq:c3-3-5}), we obtain

\[
\begin{aligned}{F(z)}&{{}=\lim_{k\to\infty}\int_{\Gamma^{*}}e^{2\pi i z\cdot t}f_{k}(t)\;d t}\\ {}&{{}=\lim_{k\to\infty}\;\int_{\Gamma^{*}}e^{2\pi i z\cdot t}\Biggl\{\int_{|\xi|\leq k}F(\xi)e^{-2\pi i t\cdot\xi}\;d\xi\Biggr\}\;d t}\\ {}&{{}=\lim_{k\to\infty}\;\int_{|\xi|\leq k}F(\xi)\Biggl\{\int_{\Gamma^{*}}e^{2\pi i(z-\xi)\cdot t}\;d t\Biggr\}\;d\xi}\\ {}&{{}=\int_{E_{n}}F(\xi)K(z-\xi)\;d\xi,}\\ \end{aligned}
\]

and the theorem is proved.

In the classical one-dimensional case the Poisson kernel\index{Poisson kernel@Poisson kernel!for tube domain@for tube domain}

\[
P(x,y)=\frac{1}{\pi}\frac{y}{x^{2}+y^{2}}
\]

can be expressed in terms of the Cauchy kernel associated with the upper half-plane in the following simple way: For $z = x + iy$, y > 0,

\[
P(x,y)=|K(z)|^{2}/K(2iy).
\]

It is natural, therefore, to extend this definition whenever the Cauchy kernel is defined. Thus, if $T_{\Gamma}$ is a tube whose base $\Gamma$ is a regular cone and K is the corresponding Cauchy kernel we define the Poisson kernel associated with $T_{\Gamma}$ by putting

\[
\mathcal{P}(x,y)=|K(x+i y)|^{2}/K(2i y)
\]

whenever $z = x + iy \in \Gamma$.

We have shown that $K(x + iy)$, considered as a function of $x$, belongs to $L^{2}(E_{n})$; hence, $\mathcal{P}(\cdot, y)$ must belong to $L^{1}(E_{n})$ for each $y$ in $\Gamma$. Moreover, it is not hard to show that $\mathcal{P}(\cdot, y)$ belongs to $L^{\infty}(E_{n})$. In order to do this it clearly suffices to show that $K(x + iy)$ is bounded independently of $x$ for each $y$ in $\Gamma$. But

\[
\left|K(x+i y)\right|=\left|\int_{\Gamma^{*}}e^{2\pi i(x+i y)\cdot t}d t\right|\leq\int_{\Gamma^{*}}e^{-2\tau y\cdot t}d t=K(i y).
\]

Thus, it suffices to show that $K(iy)$ is finite, if $y \in \Gamma$. In order to do this we first claim that, if $y \in \Gamma$, then there exists a $\delta = \delta_y > 0$ such that $\delta|t| \leq y \cdot t$ for all $t$ in $\Gamma^*$. It is enough to show this when $t \in \Gamma^*$ satisfies

$|t| = 1$. From the definition of $\Gamma^*$ we have $0 \leq y \cdot t$. On the other hand, equality is not possible; otherwise, since $\Gamma$ is open, we could find a $u \in E_n$ of small enough absolute value to guarantee that $y + u \in \Gamma$ and such that $(y + u) \cdot t = u \cdot t < 0$, contradicting the fact that $t \in \Gamma^*$. Since the intersection of $\Gamma^*$ with the surface, $\Sigma$, of the unit sphere of $E_n$ is compact, the existence of $\delta_y > 0$ now follows from the fact that $0 < y \cdot t$ for all $t \in \Gamma^* \cap \Sigma$. Consequently,

\[
\intop_{\Gamma^{*}}e^{-2\pi y\cdot t}\,dt\;\leq\intop_{\Gamma^{*}}e^{-2\pi\delta|t|}\,dt\;<\infty
\]

for $y \in \Gamma$. Since $L^q(E_n) \supset L^1(E_n) \cap L^\infty(E_n)$ for $1 \leq q \leq \infty$,\footnote{If $f \in L^1 \cap L^\infty$ and $1 < q < \infty$ then $\int |f|^q = \int |f|^{q-1}|f| \leq \|f\|_\infty^{q-1}\|f\|_1 < \infty$.} we have the result.

\begin{corollary}
\label{cor:c3-3-7}
For each $y \in \Gamma$, $\mathcal{P}(\cdot, y) \in L^q(E_n)$ whenever $1 \leq q \leq \infty$. It follows that if $f \in L^p(E_n)$, $1 \leq p \leq \infty$, then
\[
u(x+i y)=\int_{E_{n}}f(x-t)\mathcal{P}(t,y)d t
\]

is defined for all $z = x + iy \in T_{\Gamma}$. As was the case with the Poisson integrals introduced in Chapter I, it can be shown that $u(x + iy)$ tends to $f(x)$ in the $L^{p}$ norm. That is,
\begin{equation}\label{eq:c3-3-8}
\lim_{y\to0,y\in\Gamma}\int\limits_{E_{n}}|u(x+iy)-f(x)|^{p} dx=0.
\end{equation}

\end{corollary}

This is an immediate consequence of the fact that the kernel $\mathcal{P}$ is an approximation of the identity (see the footnote following the end of the proof of Theorem \ref{thm:c2-2-1} of Chapter II and the last part of the proof of Theorem \ref{thm:c3-5-6} below). We mean by this that the following three properties are satisfied by $\mathcal{P}$:

(i) $\mathcal{P}(x, y) \geq 0$;

(ii) $\int_{E_n} \mathcal{P}(x, y) \, dx = 1$ for all $y \in \Gamma$;

(iii) if $\delta > 0$ then $\int_{|x| > \delta} \mathcal{P}(x, y) \, dx \to 0$ as $y \in \Gamma$ tends to 0.

Property (i) is obvious. The second property follows from \eqref{eq:c3-3-5} after dividing both sides by $K(2iy)$. In order to prove (iii) it suffices to find a function $\psi$ such that

(a) $\psi$ is continuous on $E_{n}$,

\[
\begin{array}{r}{\mathrm{(b)}\lim_{u\in\Gamma,\;y\rightarrow0}\int_{E_{n}}\mathcal{P}(x,y)\psi(x)\;d x=1,}\end{array}
\]

(c) $|\psi(x)| < 1$ when $x \neq 0$ and $\psi(x) \to 0$ as $|x| \to \infty$.

For, if such a function exists, then, by (b),

\[
1=\lim_{y\in\Gamma, y\to0}\Bigg\{\int\limits_{|x|\leq\delta}\psi(x)\mathcal{P}(x,y)\;{d}x\;+\;\int\limits_{|x|>\delta}\psi(x)\mathcal{P}(x,y)\;{d}x\Bigg\}.
\]

From (a) and (c) we know that there exists $\varepsilon > 0$ such that $|\psi(x)| \leq 1 - \varepsilon$ if $|x| > \delta$. Thus (using (i) and (ii), as well),

\[
\begin{aligned}1\leq&\lim_{y\in\Gamma,y\to0}\Biggl\{\int_{|x|\leq\delta}\mathcal{P}(x,y)dx+(1-\varepsilon)\int_{|x|>\delta}\mathcal{P}(x,y)dx\Biggr\}\\=&\lim_{y\in\Gamma,y\to0}\Biggl\{\int_{\widetilde{E}_{n}}\mathcal{P}(x,y)dx-\varepsilon\int_{|x|>\delta}\mathcal{P}(x,y)dx\Biggr\}\\=&\lim_{y\in\Gamma,y\to0}\Biggl\{1-\varepsilon\int_{|x|>\delta}\mathcal{P}(x,y)dx\Biggr\}.\end{aligned}
\]

But this clearly implies (iii).

The existence of such a function $\psi$ will be an easy consequence of the following representation theorem:

\begin{theorem}
\label{thm:c3-3-9}
If $F \in H^2(T_\Gamma)$ then
\[
F(z)=\int\limits_{E_{n}}\mathcal{P}(x-t,y)F(t)dt
\]

for all $z = x + iy$ in $T_{\Gamma}$.
\end{theorem}

\textsc{Proof.}\ Let $w = u + iv$ be a point of $T_{\Gamma}$. Then, for $z \in T_{\Gamma}$,

\[
|K(z+w)|=\left|\int_{\Gamma^{\bullet}}e^{2\pi i(x+u)\cdot t}e^{-2\pi(y+v)\cdot t}d t\right|\leq\int_{\Gamma^{\bullet}}e^{-2\pi v\cdot t}d t=M_{v}<\infty.
\]

Thus, $F(z)K(z + w)$, as a function of $z$, belongs to $H^2(T_r)$ and has norm not exceeding $\|F\|M_v$. We can thus apply Theorem \ref{thm:c3-3-6} to this function and obtain

\begin{equation}\label{eq:c3-3-10}
F(z)K(z+w)=\int\limits_{E_{n}}K(z-t)F(t)K(t+w)d t,
\end{equation}

for all $z \in T_{\Gamma}$

If we let $w = -x + iy$ then $K(z-t)K(t+w) = |K(z-t)|^{2}$ and $K(z+w) = K(2iy)$; therefore \eqref{eq:c3-3-10} in this case is equivalent to

\[
F(z)=\int\limits_{E_{n}}F(t)\{|K(z-t)|^{2}/K(2i y)\}d t=\int\limits_{E_{n}}F(t)\mathcal{P}(x-t,y)d t,
\]

which is the result we are seeking.

We now pass to the construction of the function $\psi$. Choose $\varphi \geq 0$ that is continuous, has compact support contained in $\Gamma^*$ and such that $\int_{E_n} \varphi(t) \, dt = 1$ (this can be done since $\Gamma$ is regular). We claim that

\[
\psi(x)=\int_{E_{n}}e^{2\pi i x\cdot t}\varphi(t)\;d t=\int_{\Gamma^{*}}e^{2\pi i x\cdot t}\varphi(t)\;d t
\]

satisfies (a), (b), and (c). Since $\varphi \in L^{1}(E_{n})$, property (a) is certainly true (see Theorem \ref{thm:c1-1-1}, part (b), in the first chapter). The fact that $\psi(x)$ tends to 0 as $|x|$ tends to $\infty$ is a special case of the Riemann-Lebesgue theorem. If $|\psi(x)| = 1$, say $\psi(x) = e^{2\pi i\theta}$, then

\[
1=\int\limits_{E_{n}}e^{2\pi\mathrm{i}\left[(x\cdot t)-\theta\right]}\varphi(t)dt=\int\limits_{E_{n}}\varphi(t)\cos2\pi\left[(x\cdot t)-\theta\right]dt
\]

(the imaginary part of the first integral must vanish since 1 is a real number). If $x \neq 0$ then $\cos 2\pi[(x \cdot t) - \theta]$, as a function of $t$, must be strictly less than 1 in a subset of the support of $\varphi$ having positive measure. This, together with the assumptions that $\varphi$ is nonnegative and

\[
\int\limits_{E_{n}}\varphi(t)dt=1,
\]

implies $\left|\psi(x)\right| < 1$ if $x \neq 0$. In order to show that (b) holds we consider

\[
F(z)=\int\limits_{\Gamma^{\bullet}}e^{2\pi i z\cdot t}\varphi(t)\,dt
\]

for $z \in T_{\Gamma}$. By Theorem \ref{thm:c3-3-1} $F$ belongs to $H^2(T_{\Gamma})$ and it is clear that

\[
F(x)=\lim_{y\in\Gamma,y\to0}F(x+iy)=\int\limits_{\Gamma^{*}}e^{2\pi i x\cdot t}\varphi(t)dt=\psi(x)
\]

in the $L^{2}$-norm; moreover, by the Lebesgue dominated convergence theorem,

\begin{equation}\label{eq:c3-3-11}
F(0)=\lim_{y\in\Gamma,y\to0}F(0+iy)=\psi(0)=1.
\end{equation}

Applying Theorem \ref{thm:c3-3-9} to this function F we have

\begin{equation}\label{eq:c3-3-12}
F(x+i y)=\int_{E_{n}}\mathcal{P}(x-t,y)F(t)d t=\int_{E_{n}}\mathcal{P}(x-t,y)\psi(t)d t.
\end{equation}

Since $\mathcal{P}(-t,y) = \mathcal{P}(t,y)$, \eqref{eq:c3-3-12} (with x = 0) together with \eqref{eq:c3-3-11} gives us property (b).

\section{The Paley-Wiener Theorem}

This section will be devoted to an application of the theory of $H^{2}$ spaces. We shall show how the classical Paley--Wiener theorem can be extended to n-dimensions and proved by using some of the results we have established in this chapter. We begin by discussing the one-dimensional case.

Suppose $F$ is an entire function on the complex plane $C_1$. We say that $F$ is of exponential type\index{exponential type@exponential type} $\sigma > 0$ if for each $\varepsilon > 0$ there exists a constant $A_\varepsilon$ such that

\[
\left|F(z)\right|\leq A_{\varepsilon}e^{(\sigma+\varepsilon)|z|}
\]

for all $z \in C_1$. An example of such a function can be obtained in the following way: Suppose $f \in L^2(-\tau, \tau)$, we then define an entire function, $F$, by letting

\[
F(z)=\int\limits_{-\tau}^{\tau}f(t)e^{2\pi i z t}d t
\]

for $z \in \mathbb{C}_1$. Then

\[
|F(z)|=|F(x+i y)|\leq\sqrt{2\tau}\left(\intop_{-\tau}^{\tau}|f(t)|^{2}d t\right)^{1/2}e^{2\pi\tau|y|}\leq A e^{2\pi\tau|z|}.
\]

Hence, $F$ is of exponential type $\sigma = 2\pi\tau$. The Paley-Wiener theorem asserts that the converse of this is true, if we know that $F$ restricted to the real axis belongs to $L^{2}(-\infty, \infty)$. More precisely:

\begin{theorem}
\label{thm:c3-4-1}
Suppose $F \in L^2(-\infty, \infty)$. Then $F$ is the Fourier transform of a function vanishing outside $[-(\sigma/2\pi), \sigma/2\pi] = [-\tau, \tau]$ if and only if $F$ is the restriction to the real axis of an entire function of exponential type $\sigma$.
\end{theorem}

We have just shown the ``only if'' part of this theorem. \footnote{According to the notation established in the first chapter, we have shown that if $F$ is the inverse Fourier transform of a function $f$ vanishing outside the interval $[-\tau, \tau]$ then it is the restriction to the real axis of a function of exponential type $\sigma$. This is obviously equivalent to having $F$ the Fourier transform of the function $g$ whose values are given by $g(t) = f(-t)$.} The ``if'' part will be established with the help of the following Phragmén-Lindelöf\index{Phragmen-Lindelof theorems@Phragmen-Lindelof theorems} type result (compare with 5.2 in Chapter II):

\begin{lemma}
\label{lem:c3-4-2}
Let $S$ be the region in $C_1$ bounded by two rays meeting at the origin at an angle $\pi/\alpha$. Suppose $f$ is analytic on $\bar{S}$ and satisfies $|f(z)| \leq A \exp \{|z|^{\beta}\}$ for $0 \leq \beta < \alpha$ and $z \in S$. Then, the condition $|f(z)| \leq M$ on the two bounding rays implies $|f(z)| \leq M$ for all $z \in S$.
\end{lemma}

\textsc{Proof.}\ By rotating, if necessary, we can assume that the rays form angles $\pi/2\alpha$ and $-\pi/2\alpha$ with the real axis. Let $F(z) = f(z)\exp\{-\varepsilon z^{\prime}\}$ where $\beta < \gamma < \alpha$ and $\varepsilon > 0$. It follows that $|F(z)| \leq |f(z)|$ $\leq M$ on the two bounding rays. Moreover, on the arc $R = |z| = |re^{i\theta}|$, $-(2\pi/2\alpha) \leq \theta \leq \pi/2\alpha$, $|F(z)| \leq A \exp\{R^{\beta} - \varepsilon R^{\gamma} \cos \gamma\pi/2\alpha\}$. But the last expression tends to 0 as $R \to \infty$. Thus, $|F(z)| \leq M$ on this arc provided $R$ is large enough. The maximum-modulus principle, thus, implies $|F(z)| \leq M$ for all $z \in \overline{S}$ with $|z| \leq R$. Since $R$ can be chosen arbitrarily large this shows that $|F(z)| \leq M$ for all $z \in \overline{S}$. Hence, $|f(z)| \leq M \exp\{\varepsilon r^{\gamma} \cos \gamma\theta\}$ for all $z = re^{i\theta} \in \overline{S}$. Letting $\varepsilon \to 0$ we obtain our result.

\begin{lemma}
\label{lem:c3-4-3}
Suppose $F$ is of exponential type $\sigma$ and $|F(x)| \leq 1$ for $x$ real then $|F(x + iy)| \leq \exp\{\sigma|y|\}$ for all complex numbers $z = x + iy$.
\end{lemma}

\textsc{Proof.}\ For $\varepsilon > 0$ set $F_{\varepsilon}(z) = F(z)e^{i(\sigma + \varepsilon)z}$. Since $F$ is of exponential type $\sigma$

\[
|F_{\varepsilon}(i y)|=|F(i y)|e^{-(\sigma+\varepsilon)y}\leq A_{\varepsilon},
\]

for all nonnegative $y$. We also have $|F_\varepsilon(x)| \leq 1$ for all real $x$. This gives us a bound for $F$ on the positive $x$ and $y$ axes. Moreover, we certainly can find $B$ so that

\[
\big|F_{\varepsilon}(z)\big|\leq A_{\varepsilon}e^{(\sigma+\varepsilon)(|z|-y)}\leq A_{\varepsilon}e^{2(\sigma+\varepsilon)|z|}\leq B e^{|z|^{3:2}}
\]

We can therefore apply Lemma \ref{lem:c3-4-2} with $\beta = \frac{3}{2} < 2 = \alpha$ and obtain

\[
\left|F_{\varepsilon}(z)\right|\leq\max\left\{A_{\varepsilon},1\right\}=A
\]

for all $z = x + iy$ such that $x \geq 0$ and $y \geq 0$. If we now repeat this argument for the second quadrant, we can then apply Lemma \ref{lem:c3-4-2} to $F_\varepsilon$ restricted to the upper half-plane and $\beta = 0 < 1 = \alpha$ to obtain $|F_\varepsilon(x + iy)| \leq 1$ for $y \geq 0$. Letting $\varepsilon$ tend to zero we obtain $|F(z)| = |F(x + iy)| \leq \exp\{\sigma y\}$ for $y \geq 0$. The lemma is then established by applying this result to $G(z) = F(-z)$.

\begin{lemma}
\label{lem:c3-4-4}
Suppose F is of exponential type $\sigma$ and that its restriction to the x axis has $L^{2}$-norm not exceeding 1. Then
\[
\left(\intop_{-\infty}^{\infty}|F(x+i y)|^{2}d x\right)^{1/2}\leq e^{\sigma|y|}
\]

for all real y.
\end{lemma}

\textsc{Proof.}\ Let $\varphi$ be a bounded function of a real variable with compact support and $\|\varphi\|_2 \leq 1$. We can then define $G(z) = \int_{-\infty}^{\infty} F(z + t)\varphi(t) \, dt$.

G is, then, analytic and for $\varepsilon > 0\ |G(z)|\leq\int_{-\infty}^{\infty}A_{\varepsilon}e^{(\sigma+\varepsilon)|z|}e^{(\sigma+\varepsilon)|t|}|\varphi(t)|\,dt=B_{\varepsilon}e^{(\sigma+\varepsilon)|z|}$; thus, $G$ is also of exponential type $\sigma$. Furthermore, Schwarz's inequality implies that

\[
|G(x)|\leq\left(\int_{-\infty}^{\infty}|F(x)|^{2}d x\right)^{1/2}\left(\int_{-\infty}^{\infty}|\varphi(x)|^{2}d x\right)^{1/2}\leq1.
\]

Consequently, we can apply Lemma \ref{lem:c3-4-3} to the function G and obtain

\begin{equation}\label{eq:c3-4-5}
\left|\int\limits_{-\infty}^{\infty}F(z+t)\varphi(t)dt\right|=\left|G(z)\right|=\left|G(x+iy)\right|\leq e^{\sigma|y|}
\end{equation}

for all real y. Thus, making use of the converse of Schwarz's inequality, taking the supremum over all such $\varphi$ in \eqref{eq:c3-4-5} we obtain Lemma \ref{lem:c3-4-4}.

We can now finish the proof of the Paley--Wiener theorem. Suppose $F$ is of exponential type $\sigma$ and its restriction to the real axis belongs to $L^2(-\infty, \infty)$. We shall show that the inverse Fourier transform of this restriction vanishes (a.e.) outside the interval $[-(\sigma/2\pi), \sigma/2\pi] = [-\tau, \tau]$. There is no loss of generality if we assume that $\int_{-\infty}^{\infty} |F(x)|^2 \, dx \leq 1$.

Let $G_{+}(z) = e^{i\sigma z} F(z)$. Then, using Lemma \ref{lem:c3-4-4},

\[
\begin{aligned}\biggl(\int\limits_{-\infty}^{\infty}|G_{+}(x+iy)|^{2}dx\biggr)^{1/2}&=e^{-\sigma y}\biggl(\int\limits_{-\infty}^{\infty}|F(x+iy)|^{2}dx\biggr)^{1/2}\\&\leq e^{-\sigma y}e^{\sigma y}=1\end{aligned}
\]

for $y \geq 0$. Thus, $G_+$ belongs to $H^2(T_R^+) = H^2(E_2^+)$, the $H^2$-space associated with the upper half-plane. It follows from the results of the last section (see, in particular, Theorem \ref{thm:c3-3-1} and Corollary \ref{cor:c3-3-4}) that there exists $g \in L^2(-\infty, \infty)$, vanishing on the negative axis, such that

\[
G_{+}(x+i y)=G_{+}(z)=\int\limits_{0}^{\infty}g(t)e^{2\pi i z t}d t
\]

for all $y > 0$. Letting $f(s) = g(\tau - s) = g[(\sigma/2\pi) - s]$, this equality is equivalent to

\[
F(z)=F(x+i y)=\int\limits_{-\infty}^{\tau}f(s)e^{-2\pi i z s}d s
\]

for all $y > 0$. Passing to the limit, $y = 0$, we see that the inverse Fourier transform of $F(x)$ vanishes for almost every $s \geq \tau$. Applying the same argument to $F(-z)$ we see that it must vanish for almost every $s \leq -\tau$ as well, and Theorem \ref{thm:c3-4-1} is proved.

We shall now extend this result to n-dimensions. In order to do this we need first to find an appropriate generalization of the notion of exponential type. Toward this end we make the following definitions.

If $\|\cdot\|$ and $|\cdot|$ are two norms on a vector space (over the real or complex numbers) we say that they are equivalent if there exist constants $c_1$ and $c_2$ such that $0 < c_1 \leq \|x\|/|x| \leq c_2 < \infty$ for all nonzero vectors $x$. It is well known that any norm $\|\cdot\|$ on $E_n$ is equivalent to the Euclidean norm $|\cdot|$. The unit sphere with respect to such a norm, $\|\cdot\|$, is the set $K = \{x \in E_n; \|x\| \leq 1\}$. It is clear that $K$ is convex, compact and symmetric (that is, $x \in K$ implies $-x \in K$). If such a subset of $E_n$ has a non-empty interior it is called a symmetric body. It is not hard to show that $K \subset E_n$ is a symmetric body if and only if it is the unit sphere with respect to a norm that is equivalent to the Euclidean norm.\index{equivalent norms@equivalent norms!in En@in $E_n$}\index{symmetric body@symmetric body}

Suppose $K$ is any subset of $E_n$, then the set $K^* = \{y \in E_n; x \cdot y \leq 1$ for all $x \in K\}$ is called the polar set of $K$. For example, if, for $p \geq 1$, $K = \{x = (x_1, x_2) \in E_2; |x_1|^p + |x_2|^p \leq 1\}$ it is easy to check that $K^* = \{y = (y_1, y_2) \in E_2; |y_1|^q + |y_2|^q \leq 1\}$, where $1 = 1/p + 1/q$. It is also clear that if $K$ is a symmetric body so is $K^*$. Thus, $K^*$ is the unit sphere with respect to a norm $\|\cdot\|^*$, called the dual norm\index{norm@norm!dual@dual} to the one for which $K$ is the unit sphere. We also have the equivalent way of introducing the dual norm:\index{polar set@polar set}

\begin{equation}\label{eq:c3-4-6}
\|y\|^{*}=\sup_{x\in K}|x\cdot y|.
\end{equation}

\begin{lemma}
\label{lem:c3-4-7}
Suppose $K \subset E_n$ is convex, closed and that $0 \in K$. Then $K^{**} = (K^*)^* = K$.
\end{lemma}

\textsc{Proof.}\ It is obvious that $K \subset K^{**}$. Thus, it suffices to show that if $x_0 \notin K$ then $x_0 \notin K^{**}$. Given such a point $x_0$ choose $y \in K$ so that $|y - x_0|$ is minimal. We now partition $E_n$ into the two subsets $\{x \in E_n; x \cdot (x_0 - y) > y \cdot (x_0 - y)\}$ and $\{x \in E_n; x \cdot (x_0 - y) \leq y \cdot (x_0 - y)\}$. Since $x_0 \cdot (x_0 - y) - y \cdot (x_0 - y) = |x_0 - y|^2 > 0$ it is clear that $x_0$ belongs to the former. We claim that $K$ is contained in the latter. If this were not the case, there would exist $y_1 \in K$ such that $(y_1 - y) \cdot (x_0 - y) > 0$. Let us choose $\alpha < 1$ so that $0 < \alpha < 2(y_1 - y) \cdot (x_0 - y)/|y_1 - y|^2$. Since $K$ is convex, $w = (1 - \alpha)y + \alpha y_1 \in K$ and we have

\[
\begin{aligned}\left|w-x_{0}\right|^{2}=\alpha\{\alpha|y_{1}-y|^{2}-2(y_{1}-y)\cdot(x_{0}-y)\}\\ +|y-x_{0}|^{2}<|y-x_{0}|^{2},\end{aligned}
\]

contradicting the fact that $|y - x_{0}|$ is minimal.

Since $0 \in K$ it follows that $y \cdot (x_0 - y) \geq 0$. Thus, we can find a positive constant, $\varepsilon$, such that $x_0 \cdot (x_0 - y) > \varepsilon$ while $x \cdot (x_0 - y) \leq \varepsilon$ for all $x \in K$ (if $y \cdot (x_0 - y) > 0$ we can choose $\varepsilon = y \cdot (x_0 - y)$); while any positive

$\varepsilon < x_0 \cdot (x_0 - y)$ can be chosen if $y \cdot (x_0 - y) = 0$. Letting $v = (x_0 - y)/\varepsilon$ this means

\[
K\subset\{x\in E_{n};x\cdot v\leq1\}\quad\text{and}\quad x_{0}\cdot v>1.
\]

But this means that $v \in K^*$ and that $x_0$ cannot belong to $K^{**}$; thus, the lemma is proved.

Applying this result and the relation \eqref{eq:c3-4-6} we therefore obtain

\begin{equation}\label{eq:c3-4-8}
\|x\|=\|x\|^{**}=\sup_{y\in K^{*}}|x\cdot y|.
\end{equation}

If $z = (z_1, z_2, \ldots, z_n) \in \mathbb{C}_n$ it is then natural to define $\|z\|$ by extending \eqref{eq:c3-4-8}:

\begin{equation}
\tag{4.8$'$}\label{eq:c3-4-8p}
\|z\|=\sup_{y\in K^{*}}|z\cdot y|=\sup_{y\in K^{*}}|z_{1}y_{1}+\cdots+z_{n} y_{n}|.
\end{equation}

We now say that an entire function, $F$, on $C_n$ is of exponential type\index{exponential type@exponential type!type K@type $K$} $K$, where $K$ is a symmetric body, if for each $\varepsilon > 0$ there exists a constant $A_\varepsilon$ such that

\[
\left|F(z)\right|\leq A_{\varepsilon}e^{2\pi(1+\varepsilon)z\varepsilon}
\]

for all z. The class of all functions of exponential type K will be denoted by $\mathcal{E}(K)$.

We can now state an n-dimensional version of the Paley--Wiener theorem:

\begin{theorem}
\label{thm:c3-4-9}
Suppose $F \in L^2(E_n)$. Then F is the Fourier transform of a function vanishing outside a symmetric body K if and only if F is the restriction to $E_n$ of a function in $\mathcal{E}(K^*)$.
\end{theorem}

\textsc{Proof.}\ If F is the Fourier transform of a function, f, vanishing outside K then it is easy to check that

\begin{equation}\label{eq:c3-4-10}
F(z)=\int\limits_{E_{n}}e^{-2\pi i z\cdot t}f(t)d t=\int\limits_{K}e^{-2\pi i x\cdot t}e^{2\pi y\cdot t}f(t)d t
\end{equation}

extends F to a function in $\mathcal{E}(K^*)$. In fact, it follows immediately from \eqref{eq:c3-4-10} that

\[
|F(z)|=|F(x+i y)|\leq A e^{2\pi i|y_{i}|^{*}}.
\]

The converse will be deduced from the following n-dimensional extension of Lemma \ref{lem:c3-4-4}:

\begin{lemma}
\label{lem:c3-4-11}
Suppose $F \in \mathcal{E}(K^*)$ then
\[
\biggl(\int_{E_{n}}|F(x+i y)|^{2}d x\biggr)^{1/2}\leq e^{2\pi\|y\|^{*}}\biggl(\int_{E_{n}}|F(x)|^{2}d x\biggr)^{1/2}.
\]

\end{lemma}

We shall reduce this inequality to the one-dimensional case. Let us fix $y \neq 0$ in $E_n$ and let $e_1$ be a unit vector in $E_n$ along the direction of $y$. We now choose an orthonormal basis $\{e_1, e_2, \ldots, e_n\}$ of $E_n$. We fix $(n-1)$ real numbers $u_2, \ldots, u_n$, put $\alpha = \sum_{j=2}^{n} u_j e_j$ and define

\[
\varphi(w_{1})=F(w_{1}e_{1}+\alpha).
\]

It is clear that $\varphi$ is an entire function of the (single) complex variable $w_1 = u_1 + iv_1$. Moreover, it is easy to see that $\varphi$ is of exponential type $2\pi \|e_1\|^*$. In fact, given $\varepsilon > 0$, the fact that $F \in \mathcal{E}(K^*)$ implies the existence of a constant $A_\varepsilon$ such that

\[
\begin{aligned}{|\varphi(w_{1})|}&{{}\;\leq\;A_{\varepsilon}\exp\left\{2\pi\|w_{1}e_{1}+\alpha\|^{*}(1+\varepsilon)\right\}}\\ {}&{{}\;\leq\;[A_{\varepsilon}\exp\left\{2\pi(1+\varepsilon)\|\alpha\|^{*}\right\}]e^{2\pi\|e_{1}\|^{*}(1+\varepsilon)|w_{1}|}}\\ {}&{{}\;=\;A_{\varepsilon}^{\prime}e^{2\pi\|e_{1}\|^{*}(1+\varepsilon)|w_{1}|}.}\\ \end{aligned}
\]

Thus, by Lemma \ref{lem:c3-4-4}.

\[
\intop_{-\infty}^{\infty}|\varphi(u_{1}+i v_{1})|^{2}d u_{1}\leq e^{4\pi||e_{1}||^{*}|v_{1}|}\intop_{-\infty}^{\infty}|\varphi(u_{1})|^{2}d u_{1}
\]

\[
v_{1}\in(-\infty,\infty)
\]

Choosing $v_{1}$ so that $y = v_{1}e_{1}$ this inequality becomes

\[
\int_{-\infty}^{\infty}\left|F\left(i y+\sum_{j=1}^{n}u_{j}e_{j}\right)\right|^{2}d u_{1}\leq e^{4\pi\|y\|^{*}}\int_{-\infty}^{\infty}\left|F\left(\sum_{j=1}^{n}u_{j}e_{j}\right)\right|^{2}d u_{1}.
\]

Integrating both sides with respect to $u_{2}, u_{3}, \ldots, u_{n}$ we obtain Lemma \ref{lem:c3-4-11}.

We can now finish the proof of Theorem \ref{thm:c3-4-9}. Suppose $F$ is the restriction to $E_n$ of a function in $\mathcal{E}(K^*)$. For simplicity, we shall denote the latter also by $F$. By Lemma \ref{lem:c3-4-11}, we then have $F \in H^2(T_B)$ for all bounded bases $B$. Thus, by the general representation Theorem \ref{thm:c3-2-3} there exists $f$ such that

\[
F(z)=\int\limits_{E_{n}}e^{2\pi i z\cdot t}f(t)\,dt
\]

for all $z = x + iy \in T_B$. We can assume $0 \in B$; then Plancherel's theorem assures us that $f \in L^2(E_n)$ and $\|f\|_2^2 = \int_{E_n} |F(x)|^2 \, dx$. Thus, we see that $F$ is the inverse Fourier transform of $f$ (or, equivalently, the Fourier transform of $f(-t)$) and the theorem will be established if we show that $f$ vanishes (a.e.) outside $K$. In order to do this we first observe that the Plancherel theorem implies that

\[
\int\limits_{E_{n}}|F(x+i y)|^{2}d x=\int\limits_{E_{n}}e^{-4\pi y\cdot t}|f(t)|^{2}d t
\]

for all $y \in E_n$ (since $B$ can be chosen to be arbitrarily large). Thus, by

Lemma \ref{lem:c3-4-11}.

\begin{equation}\label{eq:c3-4-12}
\int_{E_{n}}|f(t)|^{2}e^{-4\pi y\cdot t}\;d t\;\leq\;e^{4\pi\|y\|^{\bullet}}\;\int_{E_{n}}|f(t)|^{2}\;d t
\end{equation}

for all $y \in E_n$. We claim that this inequality can hold only if $f$ vanishes almost everywhere outside $K$. Suppose $t_0 \notin K$; then, by Lemma \ref{lem:c3-4-7}, there exists $y_0 \in K^*$ such that $(t_0 \cdot y_0) < -1$ (we are also using the fact that $K^*$ is symmetric). Hence, it is clear that we can find $\delta > 0$ and a neighborhood $N = N(t_0)$ of $t_0$ such that $(t \cdot y_0) < -(1 + \delta)$ for all $t \in N$. Thus, by \eqref{eq:c3-4-12}, for $y = \rho y_0$ ($\rho > 0$)

\[
\begin{aligned}{\intop_{N}\bigl|f(t)\bigr|^{2}e^{4\pi\rho(1+\delta)}d t}&{{}\leq\intop_{N}\bigl|f(t)\bigr|^{2}e^{-4\pi y\cdot t}d t}\\ {}&{{}\leq\bigl\|f\bigr\|_{2}^{2}e^{4\pi\rho\|y_{0}\|^{*}}.}\\ \end{aligned}
\]

Since $y_0 \in K^*$ we must have $\|y_0\|^* \leq 1$. Thus, we have

\[
\left(\int\limits_{N}\left|f(t)\right|^{2}dt\right)e^{4\pi\rho(1+\delta)}\leq\|f\|_{2}^{2}e^{4\pi\rho}
\]

for all $\rho > 0$. If $\int_{N} |f(t)|^{2} dt$ were not zero, this would imply

\[
e^{4\pi\rho\delta}\;\leq\;\|f\|_{2}^{2}\bigg/\intop_{N}\lvert f(t)\rvert^{2}\;{d}t,
\]

which is clearly impossible for large $\rho$. Thus $f(t) = 0$ for almost all $t \in N$ and the theorem follows.

\section{The  $H^{p}$ Theory}

Up to this point we have considered, almost exclusively, properties of analytic functions on a tube that belong to $H^2$. In particular, we have considered only the existence of $L^2$ boundary values. In this section we shall derive some results concerning the existence of pointwise limits as we approach a point on the boundary of the base, as well as the existence of $L^p$ limits, $p > 0$.

We begin by examining the special case of the tube $T_{r}$ where $\Gamma$ is the first\index{octant@octant!first@first} octant:

\[
\Gamma=\{y=(y_{1},\ldots,y_{n})\in E_{n};y_{1}>0,\ldots,y_{n}>0\}.
\]

Γ is, then, a cone and it is clear that $\Gamma^* = \overline{\Gamma}$. Thus, the Cauchy kernel associated with the tube $T_{\Gamma}$ is

\[
K(z)=\int\limits_{0}^{\infty}\int\limits_{0}^{\infty}\cdots\int\limits_{0}^{\infty}e^{2\pi i(z_{1}t_{1}+z_{2}t_{2}+\cdots+z_{n}t_{n})}d t_{1}d t_{2}\cdots d t_{n}=\prod_{j=1}^{n}\frac{--1}{2\pi i z_{j}},
\]

where $z = (z_1, z_2, \ldots, z_n) \in T_\Gamma$. That is, $K(z)$ is simply the product of $n$ one-dimensional Cauchy kernels associated with the upper half-plane. It follows that the Poisson kernel $\mathcal{P}(x, y)$ (see 3.7 and the discussion preceding it) is the product of $n$ one-dimensional Poisson kernels:

\[
\mathcal{P}(x,y)=\prod_{j=1}^{n}\frac{1}{\pi}\frac{y_{j}}{x_{j}^{2}+y_{j}^{2}},\qquad y_{1},y_{2},\ldots,y_{n}>0.
\]

This kernel, therefore, is a special case of the iterated Poisson kernels that were studied in the third section of the previous chapter (see, in particular, Theorem \ref{thm:c2-3-22}, Corollary \ref{cor:c2-3-23} and the material preceding these results). By identifying the components $z_j = x_j + iy_j$ of a point $z = (z_1, z_2, \ldots, z_n) \in T_\Gamma$ with the ordered pairs $(x_j, y_j)$ we may consider the tube $T_\Gamma$ to be the Cartesian product $E_2^+ \times E_2^+\times\cdots\times E_2^+$ of $n$ upper half-planes of two dimensional real Euclidean space. Thus, it is consistent with the definitions introduced in the third section of Chapter II to say that a function $u$, defined in $T_\Gamma$,

(i) has a nontangential limit $l$ in each variable at $x \in E_n$ if $u(\zeta) = u(\xi + i\eta) = u(\xi_1 + i\eta_1, \ldots, \xi_n + i\eta_n)$ tends to $l$ as the point $\zeta = (\xi_1, \eta_1; \xi_2, \eta_2; \ldots; \xi_n, \eta_n)$ tends to $x = (x_1, x_2, \ldots, x_n) = (x_1, 0; x_2, 0; \ldots; x_n, 0)$ within the Cartesian product

\[
\gamma_{\alpha}(x)=\Gamma_{\alpha_{1}}(x_{1})\times\Gamma_{\alpha_{2}}(x_{2})\times\cdots\times\Gamma_{\alpha_{n}}(x_{n})\subset T_{\Gamma},
\]

for each $n$-tuple $\alpha = (\alpha_1, \alpha_2, \ldots, \alpha_n)$ of positive real numbers\footnote{We are using the notation developed in Section 3 of Chapter II: 
\[
\Gamma_{\alpha j}(x_{j})=\{(\xi_{j},\eta_{j})\in E_{2}^{+}\;;|\xi_{j}-x_{j}|<\alpha_{j}\eta_{j}\},j=1,2,\ldots,n.
\]};

(ii) is nontangentially bounded in each variable at $x \in E_n$ if $u$ is bounded in the set $\gamma_\alpha(x) \cap \{\xi + i\eta \in T_\Gamma; \eta_1, \eta_2, \ldots, \eta_n \leq 1\}$ for each $n$-tuple $\alpha$ of positive real numbers.

The existence of pointwise limits and limits in the norm as we approach the origin within the first octant can be established for all functions in $H^{p}(T_{\Gamma})$, whenever p > 0. In fact, we shall prove the following results:

\begin{theorem}
\label{thm:c3-5-1}
Suppose F belongs to $H^p(T_\Gamma), p > 0$, where $\Gamma$ is the first octant in $E_n$, then
(a) F has a nontangential limit $F(x)$ in each variable at almost every x in $E_n$. In particular,
\[
\lim_{y\in\Gamma,y\to0}F(x+iy)=F(x)
\]

exists for almost every $x \in E_n$;
(b) $\int_{E_n}|F(x + iy) - F(x)|^p \, dx$ tends to zero as $y \in \Gamma$ tends to $0 \in E_n$.
\end{theorem}

\textsc{Proof.}\ Let us fix $(\zeta_2, \ldots, \zeta_n) = (\xi_2 + i\eta_2, \ldots, \xi_n + i\eta_n), \eta_2, \ldots, \eta_n > 0$, and consider the function g of one complex variable defined by letting $g(\zeta_1) = F(\zeta_1, \zeta_2, \ldots, \zeta_n)$ whenever $\zeta_1 = \xi_1 + i\eta_1$ belongs to the upper half-plane $F_2^+$. We first show that g belongs to $H^p(E_2^+)$. Since $|F(\zeta_1, \zeta_2, \ldots, \zeta_n)|^p$ is subharmonic as a function of $\zeta_2 = \xi_2 + i\eta_2$ (see example (3) in Section 4 and \eqref{frp:c2-5-1} in Chapter II) we have, writing $w_2 = u_2 + iv_2$,

\[
\begin{aligned}|F(\zeta_{1},\zeta_{2},\cdots,\zeta_{n})|^{p}&\leq\frac{1}{\pi\eta_{2}^{2}}\int\limits_{|w_{2}-\zeta_{2}|<\eta_{2}}|F(\zeta_{1},w_{2},\cdots,\zeta_{n})|^{p}du_{2}dv_{2}\\&\leq\frac{1}{\pi\eta_{2}^{2}}\int\limits_{0}^{\frac{2\eta_{2}}{}}\int\limits_{-\infty}^{\infty}|F(\zeta_{1},w_{2},\cdots,\zeta_{n})|^{p}du_{2}dv_{2}.\end{aligned}
\]

Repeating this argument for $\zeta_{3}, \ldots, \zeta_{n}$ we obtain

\[
\begin{aligned}|F(\zeta)|^{p}\leq&\pi^{1-n}(\eta_{2}\cdots\eta_{n})^{-2}\\&\times\int\limits_{0}^{2\eta_{n}}\cdots\int\limits_{0}^{2\eta_{2}}\left(\int\limits_{E_{n-1}}|F(\zeta_{1},w_{2},\ldots,w_{n})|^{p} du_{2}\cdots du_{n}\right)dv_{2}\cdots dv_{n}.\end{aligned}
\]

Now integrating both sides of this inequality with respect to $\xi_{1}$ we have

\[
\begin{aligned}\int\limits_{-\infty}^{\infty}|g(\xi_{1}+&i\eta_{1})|^{p}d\xi_{1}\\&\leq\pi^{1-n}(\eta_{2}\cdots\eta_{n})^{-2}\int\limits_{0}^{2\eta_{n}}\cdots\int\limits_{0}^{2\eta_{2}}\left(\int\limits_{E_{n}}|F(u+iv)|^{p}du\right)dv_{2}\cdots dv_{n}\\&\leq\pi^{1-n}(\eta_{2}\cdots\eta_{n})^{-2}\|F\|_{p}^{p}(2\eta_{2}\cdots2\eta_{n})=A_{1}<\infty.\end{aligned}
\]

This shows that g is a member of $H^{p}(E_{2}^{+})$.

For $\zeta = (\zeta_1, \zeta_2, \ldots, \zeta_n)$ in $T_\Gamma$ we let $s(\zeta) = |F(\zeta)|^{p/2}$. The last inequality can then be expressed in terms of the function $s$:

\[
\intop_{-\infty}^{\infty}[s(\xi_{1}+i\eta_{1},\zeta_{2},\ldots,\zeta_{n})]^{2}d\xi_{1}=\intop_{-\infty}^{\infty}|g(\xi_{1}+i\eta_{1})|^{p}d\xi_{1}\leq A_{1}<\infty.
\]

Since $s(\zeta_1, \zeta_2, \ldots, \zeta_n)$ is a subharmonic function of $\zeta_1$, the other variables being kept fixed, this inequality together with \eqref{eq:c2-4-12} in the second chapter assures us that

\[
\begin{aligned}{s(\xi_{1}+i[\varepsilon_{1}+\eta_{1}],\zeta_{2},\ldots,\zeta_{n})}\\ {\leq\frac{1}{\pi}\intop_{-\infty}^{\infty}s(t_{1}+i\varepsilon_{1},\zeta_{2},\ldots,\zeta_{n})\frac{\eta_{1}}{(\xi_{1}-t_{1})^{2}+\eta_{1}^{2}}d t_{1}}\\ \end{aligned}
\]

for all $\epsilon_1 > 0$. Applying the same argument to $\zeta_2, \ldots, \zeta_n$ we obtain

\begin{equation}\label{eq:c3-5-2}
\begin{aligned}{s(\xi+i[\varepsilon+\eta])}&{{}=s(\xi_{1}+i[\varepsilon_{1}+\eta_{1}],\ldots,\xi_{n}+i[\varepsilon_{n}+\eta_{n}])}\\ {}&{{}\leq\intop_{E_{n}}s(t+i\varepsilon)\mathcal{P}(\xi-t,\eta)\;d t}\\ \end{aligned}
\end{equation}

for all $\varepsilon=(\varepsilon_{1},\ldots,\varepsilon_{n})$ in $\Gamma$.

The fact that $F$ belongs to $H^p(T_\Gamma)$ implies that the family $\{f_\varepsilon\}$, where $f_\varepsilon(t) = s(t + i\varepsilon)$ for $\varepsilon$ in $\Gamma$, is uniformly bounded in the $L^2$ norm (in fact $\|f_\varepsilon\|_2^2 \leq \|F\|_p^p$). Thus, by the weak compactness of the unit sphere in $L^2(E_n)$, we can find a sequence $\{\varepsilon^{(k)}\}$ tending to $0 \in E_n$ such that $\{f_\varepsilon^{(k)}\}$ converges weakly to a function $f$ in $L^2(E_n)$ as $k \to \infty$. Putting $\varepsilon = \varepsilon^{(k)}$ in \eqref{eq:c3-5-2} and then letting $k \to \infty$, we obtain

\begin{equation}\label{eq:c3-5-3}
s(\xi+i\eta)\leq\int_{E_{\eta}}f(t)\mathcal{P}(\xi-t,\eta)d t=m(\xi+i\eta).
\end{equation}

Since $m(\xi + i\eta)$ is the iterated Poisson integral of a function in $L^2(E_n)$ it is nontangentially bounded in each variable at almost every $x \in E_n$ (see the argument following Corollary \ref{cor:c2-3-23} in the second chapter). Because of \eqref{eq:c3-5-3}, the same must be true of $s(\xi + i\eta) = |F(\xi + i\eta)|^p/2$. We can therefore apply Theorem \ref{thm:c2-3-24} of Chapter II to the real and imaginary parts of $F$ and we obtain part (a) of Theorem \ref{thm:c3-5-1}.

Part (b) now follows easily. If we let $M = M^{(n)}M^{(n-1)} \cdots M^{(1)}f$ be the $n$-fold iterate of the one-dimensional maximal functions corresponding to the variables $x_1, x_2, \ldots, x_n^9$ it then follows that $m(x + iy) \leq A(Mf)(x)$ for all $x + iy \in T_\Gamma$, where $A$ is a constant independent of $y \in \Gamma$ and $f \in L^2(E_n)$ (see the argument preceding Theorem \ref{thm:c2-3-22} in Chapter II). Thus, by \eqref{eq:c3-5-3}

\begin{equation}\label{eq:c3-5-4}
|F(x+i y)|^{p}\leq\{A(M f)(x)\}^{2}.
\end{equation}

By proving part (a) we have established the almost everywhere convergence $\lim_{y \in \Gamma, y \to 0} |F(x + iy) - F(x)| = 0$; it follows from inequality \eqref{eq:c3-5-4} that this convergence is dominated almost everywhere by a constant times $\{(Mf)(x)\}^2$:

\[
|F(x+i y)-F(x)|^{p}\leq2^{p}(|F(x+i y)|^{p}+|F(x)|^{p})\leq2^{p+1}\{A(M f)(x)\}^{2}.
\]

Since $f \in L^2(E_n)$ it follows from \eqref{eq:c2-3-20}, Chapter II, that $Mf \in L^2$. Part (b) is, therefore, an immediate consequence of the Lebesgue dominated convergence theorem.

$^{\mathrm{a}}$ If $g$ is a function of $x = (x_1, x_2, \ldots, x_n)$ then, for $1 \leq j \leq n$,

\[
(M^{(j)}g)(x)=\sup_{r>0}(2r)^{-1}\int\limits_{|t|\leq r}|g(x_{1},\ldots,x_{f}-t,\ldots,x_{n})| d t.
\]

We shall now derive some consequences of this theorem. First, suppose $\Gamma$ is the interior of the convex hull of $n$ linearly independent rays meeting at the origin. By selecting vectors $a_1, a_2, \ldots, a_n$ in the direction of these rays $\Gamma$ can be described as the open cone $\{v = v_1 a_1 + \cdots + v_n a_n \in E_n; v_1 > 0, \ldots, v_n > 0\}$. Theorem \ref{thm:c3-5-1} is then valid for a function $F$ in $H^p(T_\Gamma)$. To see this consider the linear transformation $A$ mapping the standard basis vector $e_j$ onto $a_j, j = 1, 2, \ldots, n$. Then $A$ maps the first octant onto $\Gamma$. Extending $A$ linearly to $C_n$ and putting $G(x + iy) = G(z) = F(Az) = F(Ax)$ for $z$ in the tube over the first octant, we obtain

\[
\int_{E_{n}}|G(x+i y)|^{p}d x=\frac{1}{|\det A|}\int_{E_{n}}|F(u+i A y)|^{p}d u\leq\frac{1}{|\det A|}\|F\|_{p}^{p},
\]

for all y in the first octant. Thus, G satisfies the hypotheses of Theorem \ref{thm:c3-5-1}. But this clearly implies the a.e. existence of the limit

\[
F(u)=\lim_{v\in\Gamma,v\to0}F(u+.iv)
\]

as well as the corresponding limit in the $L^{p}$-norm.

More generally, we shall say that $\Gamma$ is a polygonal\index{cone@cone!polygonal@polygonal} cone if it is the interior of the convex hull of a finite number of rays meeting in the origin, among which we can find at least n that are linearly independent. A polygonal cone is clearly a finite union of cones of the type just considered in the preceding paragraph; thus, it follows that the existence of pointwise and norm limits as we approach the origin is valid as well for functions in the class $H^p$ associated with a tube whose base is a polygonal cone.

From these considerations we can easily deduce the following theorem:

\begin{theorem}
\label{thm:c3-5-5}
Suppose $F \in H^p(T_\Gamma)$, $p > 0$, where $\Gamma$ is an open convex cone in $E_n$, then if $\Gamma_1$ is any cone whose closure is contained in $\Gamma \cup \{0\}$;
(a) $\lim_{y \in \Gamma_1, y \to 0} F(x + iy) = F(x)$ exists for almost every $x \in E_n$;
\[
\begin{array}{r l}{\mathrm{(b) }\int_{E_{n}}\left|F(x+i y)-F(x)\right|^{p}d x\to 0\text{ as }y\in\Gamma_{1}\text{ tends to }0.}\end{array}
\]

\end{theorem}

\textsc{Proof.}\ It clearly suffices to show that we can find a polygonal cone $\Gamma_0$ such that $\overline{\Gamma}_1 \subset \Gamma_0 \cup \{0\} \subset \Gamma \cup \{0\}$. In order to do this we consider the intersection of $\overline{\Gamma}_1$ with the surface of the unit sphere $\Sigma_{n-1} \subset E_n$. Let $S$ denote this intersection. For each $x \in S$ we can certainly find an open polygonal cone $\Gamma_x$ such that $x \in \Gamma_x \subset \Gamma$. Since $S$ is compact, a finite number of such cones, $\Gamma_{x_1}, \ldots, \Gamma_{x_m}$, will cover $S$. Letting $\Gamma_0$ be the convex hull of the union of these $m$ cones we then clearly have the desired inclusions $\overline{\Gamma}_1 \subset \Gamma_0 \cup \{0\} \subset \Gamma \cup \{0\}$.

In the second section of this chapter (immediately following Theorem \ref{thm:c3-2-11}) we introduced the notions of restricted\index{non-tangential limit@non-tangential limit!restricted@restricted} and unrestricted limits as we approached a boundary point of the base $B$ of a tube $T_{B}$. It follows from this proof that Theorem \ref{thm:c3-5-5} has an equivalent formulation in terms of these notions:

\begin{customthm}{5.5$'$}
\label{thm:c3-5-5p}
If $F \in H^p(T_\Gamma), p > 0$, then $F(z)$ has restricted limits almost everywhere and in the $L^p$-norm as we approach 0 within $\Gamma$.
\end{customthm}

The proof also shows that whenever $F \in H^p(T_\Gamma)$ then $F$ has restricted nontangential limits for almost every $x^0$, in the sense that the limit $F(x + iy)$ exists as $x + iy \to x^0$ when $y \in \Gamma_1$ and $x + iy \in \gamma_\alpha(x^0)$, for any proper subcone $\Gamma_1$ of $\Gamma$ and any Cartesian product cone $\gamma_\alpha(x^0)$.

With the help of the Poisson kernel $\mathcal{P}(x, y)$ associated with the tube $T_r$ it can be shown that when $p \geq 1$ we obtain unrestricted limits in the norm. More precisely, we shall show that the following result is true:

\begin{theorem}
\label{thm:c3-5-6}
Suppose $\Gamma$ is a regular open convex cone in $E_n$ and that $F \in H^p(T_\Gamma)$, $1 \leq p < \infty$, then
\[
\lim_{y\in\Gamma,y\rightarrow0}\int\limits_{E_{n}}|F(x+i y)-F(x)|^{p}d x=0,
\]

where $F(x)$ is the limit function whose existence was established in the last theorem.
\end{theorem}

\textsc{Proof.}\ We first show that the Poisson kernel $\mathcal{P}(x, y)$ has the following ``semigroup property'' \footnote{It is an easy consequence of Corollary \ref{cor:c1-1-28} in Chapter I that $u(x, y_1 + y_2) = \int_{E_n} u(t, y_2)P(x - t, y_1) \, dt$ whenever u is a harmonic function in $E_{n+1}^+$ satisfying $\|u(\cdot, y)\|_p \leq c < \infty$ for all y > 0. In the present situation we have to restrict ourselves to holomorphic functions F as the kernel $\mathcal{P}$ does not in general reproduce harmonic (nor multiply harmonic) functions.}: When $F \in H^p(T_\Gamma)$ and $y_1, y_2$ belong to $\Gamma$ then

\begin{equation}\label{eq:c3-5-7}
F(x+i(y_{1}+y_{2}))=\int\limits_{E_{n}}F(t+i y_{2})\mathcal{P}(x-t,y_{1})d t.
\end{equation}

When $p = 2$ this is an immediate consequence of Theorem \ref{thm:c3-3-9}. In general, it follows from Lemma \ref{lem:c3-2-12} that (whenever $p > 0$) $G(x + iy) = F(x + i(y + y_2))$ is uniformly bounded in $T_\Gamma$. Let $\varphi$ be a continuous non-negative function with compact support in $\Gamma^*$ and satisfying $\int_{E_n} \varphi(t) \, dt = 1$ (such a $\varphi$ exists since $\Gamma$ is regular). If $\psi(z) = \int_{E_n} e^{2\pi i z \cdot t} \varphi(t) \, dt$, Theorem \ref{thm:c3-3-1} assures us that $\psi \in H^2(T_\Gamma)$; consequently, $G_\varepsilon(z) = \psi(\varepsilon z)G(z)$ belongs

to $H^{2}(T_{\Gamma})$ for all $\varepsilon > 0$. Thus, by Theorem \ref{thm:c3-3-9},

\begin{equation}\label{eq:c3-5-8}
G_{\varepsilon}(x+i y)=\int_{E_{n}}G_{\varepsilon}(t)\mathcal{P}(x-t,y)d t.
\end{equation}

Moreover, $|\psi(z)| \leq \int_{\Gamma^*} e^{-2\pi y \cdot t} \varphi(t) \, dt \leq 1$ and $\lim_{\varepsilon \to 0} \psi(e z) = 1$. These facts, together with the fact that $\mathcal{P}(\cdot, y) \in L^q(E_n)$, where $1/p + 1/q = 1$ (see 3.7), allow us to take the limits, as $\varepsilon$ approaches zero, of both sides of equality \eqref{eq:c3-5-8} and to conclude that $G(x + iy) = \int_{E_n} G(t)\mathcal{P}(x - t, y) \, dt$. But this is precisely the desired property \eqref{eq:c3-5-7} when $y = y_1$.

Part (b) of Theorem \ref{thm:c3-5-5} now assures us that, by letting $y_{2}$ tend to 0 restrictedly, we obtain

\begin{equation}\label{eq:c3-5-9}
F(x+i y)=\int\limits_{E_{n}}F(t)\mathcal{P}(x-t,y)d t
\end{equation}

for all $y \in \Gamma$.

The rest of the proof is analogous to that of Theorem \ref{thm:c2-1-10} of Chapter II and is based on properties (i), (ii), and (iii) of the kernel $\mathcal{P}$ that were established immediately following \eqref{eq:c3-3-8}. \footnote{In fact, the proof we now give can be used in order to prove equality \eqref{eq:c3-3-8}.} Letting $\omega(r)$ denote the modulus of continuity of F and using (ii), Minkowski's ordinary inequality followed by his integral inequality and (i), in that order, we have

\[
\begin{aligned}&\left(\int_{\tilde{E}_{n}}|F(x+iy)-F(x)|^{p} dx\right)^{1/p}\\&\quad=\left(\int_{\tilde{E}_{n}}\left|\int_{\tilde{E}_{n}}\{F(x-t)-F(x)\}\mathcal{P}(t,y)dt\right|^{p} dx\right)^{1/p}\\&\quad\leq\left(\int_{\tilde{E}_{n}}\left|\int\limits_{|t|\leq\delta}\{F(x-t)-F(x)\}\mathcal{P}(t,y)dt\right|^{p} dx\right)^{1/p}\\&\quad+\left(\int_{\tilde{E}_{n}}\left|\int\limits_{|t|>\delta}\{F(x-t)-F(x)\}\mathcal{P}(t,y)dt\right|^{p} dx\right)^{1/p}\\&\quad\leq\int_{|t|\leq\delta}\left(\int_{\tilde{E}_{n}}|F(x-t)-F(x)|^{p} dx\right)^{1/p}\mathcal{P}(t,y)dt\\&\quad+\int_{|t|>\delta}\left(\int_{\tilde{E}_{n}}\|F(x-t)-F(x)\|^{p} dx\right)^{1/p}\mathcal{P}(t,y)dt\\&\quad\leq\left\{\sup_{|t|\leq\delta}\omega(|t|)\right\}\int_{\tilde{E}_{n}}\mathcal{P}(t,y)dt+2\|F\|_{p}\int_{|t|>\delta}\mathcal{P}(t,y)dt.\end{aligned}
\]

The first summand equals $\sup_{|t|\leq\delta}\omega(|t|)$ since, by (ii), $\int_{E_n}\mathcal{P}(t,y)\,dt=1$; therefore, it is less than, say, $\varepsilon>0$ if $\delta$ is sufficiently small (see the proof of Theorem \eqref{thm:c1-1-18} in Chapter I). Having chosen such a $\delta$, property (iii) then asserts that the last summand tends to 0 as $y$ approaches the origin within $\Gamma$. Thus, $(\int_{E_n}|F(x+iy)-F(x)|^p\,dx)^{1/p}<\varepsilon$ if $y\in\Gamma$ is close to 0. This proves the theorem.

\section{Further Results}

\frp{6}{1} Throughout most of this chapter we considered only tubes having an open convex base $B$. In the second section we showed that if $F \in H^2(T_B)$, where $B$ is merely open and connected, then $F$ has an analytic extension, $G$, to the convex hull $T_B^c$ of the tube $T_B$ and $G \in H^2(T_B^c)$ with $\|F\|_2 = \|G\|_2$ (see Corollary \ref{cor:c3-2-4}). This fact, as well as most of the results of the second section of this chapter, is an easy consequence of the basic representation theorem (Theorem \ref{thm:c3-2-3}). There are many other reasons for considering only convex bases. It is a classical fact that whenever $F$ is holomorphic in a tube $T_B$ with an open and connected base $B$, then it has an analytic extension, $G$, to $T_B^c$ (see Theorem 9 of Chapter V in Bochner and Martin [1]). It is easy to see that $G$ assumes only the values assumed by $F$. If this were not the case and $z_0$ were in the range of $G$ but not in the range of $F$ we would have a holomorphic function $1/(F - z_0)$ in $T_B$ which, by analytic continuation, has the analytic extension $1/(G - z_0)$ in $T_B^c$. This is clearly impossible.

These facts can be used to obtain another proof of Corollary \ref{cor:c3-2-4} that extends to give us the corresponding result for other $H^p$ spaces. Suppose $F \in H^2(T_B)$ and $h \in L^2(E_n)$, where $\int_{E_n} |h(t)|^2 \, dt = 1$. We then obtain a holomorphic function $f$ on $T_B$ by letting

\[
f(z)=f(x+i y)=\int\limits_{E_{n}}F(x+t+i y)h(t)d t
\]

whenever $z \in T_B$. By Hölder's inequality, $|f(z)| \leq \|F\|_2 \|h\|_2 = \|F\|_2$. If $G$ is the analytic extension of $F$ to $T_B^c$ then $g(z) = \int_{E_n} G(z + t) h(t) \, dt$ clearly is the analytic extension of $f$ to $T_B^c$.\footnote{It is not hard to show that the last integral is defined. By first considering a subset $B_{0} \subset B$ as in Lemma \ref{lem:c3-2-12} we can assume that F and, therefore, G is bounded; then, by choosing h in $L^{1}(E_{n}) \cap L^{2}(E_{n})$ we are assured that $G(z + t)h(t)$ is integrable. This clearly suffices for our purposes.} Thus, it follows from the above remarks that $|g(z)| \leq \|F\|_2$ for all $z \in T_B^c$. Consequently,

choosing $h$ so that $g(z) = \int_{E_n} G(z + t)h(t) \, dt = (\int_{E_n} |G(z + t)|^2 \, dt)^{1/2} = (\int_{E_n} |G(t + ty)|^2 \, dt)^{1/2}$ we see that $G \in H^2(T_B^c)$ with $\|G\|_2 \leq \|F\|_2$. If $p \geq 1$ and $\|h\|_q = 1$, where $(1/p) + (1/q) = 1$, the same argument shows that the analytic extension $G$ of an $F \in H^p(T_B)$ belongs to $H^p(T_B^c)$ and $\|G\|_p = \|F\|_p$. The result actually holds for all $p > 0$, but the proof is more difficult. It makes use of the fact that if $F \in H^p(T_B)$, $p > 0$, then $\log \|F(\cdot + iy)\|_p$ is a convex function of $y \in B$.

\frp{6}{2} Another application of the $H^{2}$ theory (and, in particular, of Theorem \ref{thm:c3-2-3}) is the following:

Let $B^+$ be an open convex subset of $E_n$ such that 0 is one of its boundary points and $B^- = \{x \in E_n; -x \in B^+\}$. Suppose that there exist two functions $F^+ \in H^2(T_B^+)$ and $F^- \in H^2(T_B^-)$ such that

\[
F^{+}(x+i0)=\lim_{y\in B^{+},y\to0}F^{+}(x+iy)=\lim_{y\in B^{-},y\to0}F^{-}(x+iy)=F^{-}(x-i0)
\]

a.e. (these limits are to be interpreted in the sense described in \eqref{eq:c3-2-8}). Then there exists an analytic function, $F$, defined on the tube over the convex hull of $B^+ \cup B^-$, such that $F(z) = F^+(z)$ if $z \in T_B^+$ and $F(z) = F^-(z)$ if $z \in T_B^-$.

To show this we first note that Theorem \ref{thm:c3-2-3} implies that there exist two functions, $f^+$ and $f^-$, such that $F^+(z) = \int_{E_n} e^{2\pi iz \cdot t} f^+(t) \, dt$ and $F^-(z) = \int_{E_n} e^{2\pi iz \cdot t} f^-(t) \, dt$. It follows that $f^+$ and $f^-$ are the Fourier transforms of $F^+(x + i0) \quad (= F^-(x - i0))$ (see the discussion connected with \eqref{eq:c3-2-8}). Consequently, $f^+(t) = f^-(t)$ a.e. Letting $f(t) = f^\pm(t)$ a.e. and using the fact that $S = \{y \in E_n; \int_{E_n} |f(t)|^2 e^{-4\pi y \cdot t} \, dt \leq M^2 < \infty\}$ is convex (see the proof of Corollary \ref{cor:c3-2-4}) we see that $F(z) = \int_{E_n} e^{2\pi iz \cdot t} f(t) \, dt$ is analytic on $T_S$, a set containing the tube over the convex hull of $B^+ \cup B^-$ (see \eqref{eq:c3-2-2}), and coincides with $F^+(z)$ on $T_B^+$ and with $F^-(z)$ on $T_B^-$.

For related results and various extensions see Streater and Wightman [1].

\frp{6}{3} Theorem \ref{thm:c3-2-11} can be extended to three dimensions. The result can be stated in the following way (see E. M. Stein, G. Weiss, and M. Weiss [1]):

Suppose $y_0$ belongs to the boundary $\partial B$ of an open convex set $B \subset E_3$. Then $y_0$ is a point of unrestricted $L^2$-convergence (that is, $\lim\limits_{y \to y_0, y \in B} F(x + iy)$ exists in the $L^2$-sense for all $F \in H^2(T_B)$) if and only if there exists a finite set $V = \{y_1, y_2, \ldots, y_n\}$ of ``vertices'' lying in $\overline{B}$ and a neighborhood $N$ of

$y_0$, so that, whenever $y \in N \cap \partial B$ and $\pi_y$ is a plane of support at $y$, then $\pi_y$ contains at least one element of $V$.

This condition is not a local one. For example, if $B$ is a circular cone this result shows that every point of $\partial B$ is a point of unrestricted $L^{2}$-convergence; however, truncating the cone by removing a neighborhood of the vertex, the remaining boundary points are no longer points of unrestricted convergence.

It is an open problem to find how this result generalizes to higher dimensions.

\frp{6}{4} If $F$ is holomorphic in the upper half-plane $E_2^+ = \{z = x + iy; y > 0\} \subset C_1$, then, except for a set of points in the $x$ axis of linear Lebesgue measure zero, the following is true: Either $F$ has a nontangential limit at $(x_0, 0)$ or $F$ maps each triangular region $\Delta(x_0; \alpha, \beta) = \{x + iy \in E_2^+\}$; $|x - x_0| < \alpha y < \beta\}$ densely into the complex plane. This is a special case of a classical result of Plessner [1] and it can be easily derived from the special case $n = 1$ of Theorem \ref{thm:c2-3-19} of Chapter II. Similarly, Theorem \ref{thm:c2-3-24} of Chapter II gives us an extension to $n$-dimensions of Plessner's result (see Calderón [4]): If $F$ is holomorphic in $T_r$, where $\Gamma$ is the first octant, then except for a subset of $E_n$ of measure zero either $F$ has a nontangential limit in each variable at $x_0 = (x_1^0, x_2^0, \ldots, x_n^0)$ or it maps densely into the complex plane each Cartesian product $\Delta(x_0; \alpha, \beta) = \Delta(x_1^0; \alpha_1, \beta_1) \times \cdots \times \Delta(x_n^0; \alpha_n, \beta_n)$, where $\alpha = (\alpha_1, \ldots, \alpha_n)$ and $\beta = (\beta_1, \ldots, \beta_n)$ belong to $\Gamma$. From this, it follows that if the real part of a holomorphic function $F = u + iv$, defined on $T_r$, has a nontangential limit in each variable at the points of a set $S \subset E_n$ then the imaginary part of $F$ has such limits at almost every point of $S$. We see this by observing that if $u$ has a nontangential limit in each variable at $x$ then $F$ cannot map $\Delta(x; \alpha, \beta)$ into a dense subset of $C_1$.

\frp{6}{5} Other notions of restricted convergence have been studied. For simplicity, suppose $u$ is a function defined in $T_{\Gamma}$, where $\Gamma$ is the first quadrant in $E_2$. We shall say that $u$ has a restricted curvilinear limit $u(x_0)$ at $x_0 \in E_2$ if $u(x_0 + iy(t)) \to u(x_0)$ as $t \to 0$, whenever the curve $y(t) = (h(t), k(t))$, $t \geq 0$, is continuous, the functions $h$ and $k$ are strictly increasing to $\infty$ and $h(0) = 0 = k(0).\footnote{More generally, we can consider nontangential restricted curvilinear limits in each variable by letting $x + iy(t)$ tend to $x_0$ within the regions $\gamma_a(x_0)$ introduced at the beginning of Section 5. The reader should not confuse this type of nontangential (or tangential) approach to $x_0$ (within $T_\Gamma$) with the nontangential (or tangential) approach of $y = y(t)$ to 0 (within $\Gamma$). We have also sacrificed generality for the sake of simplicity by considering only the two-dimensional case; these notions have an obvious extension to $n$-dimensions.}$ We have already indicated (see \eqref{frp:c2-5-7} in Chapter II) that an iterated Poisson integral $u(x + iy)$ of an $L^1$-function,

$f$, need not have unrestricted limits at almost every $x \in E_2$. One can show, however, that $u(x + iy)$ does converge almost everywhere to $f(x)$ in the restricted curvilinear sense we have just described. This can be done by appropriately modifying the argument given in Zygmund [1], Volume 2. We can extend this notion of restricted convergence to tubes having a general cone for a base and obtain another version of Theorem \ref{thm:c3-5-5p}. (In this connection see the comment immediately following Theorem \ref{thm:c3-5-5p}) The counterexample that was constructed in Section 2 shows that the situation is considerably more complicated when the base of the tube is not a cone. A further illustration of the aptness of restricted convergence when the base is a cone is provided by the examples described in 6.8 below.

\frp{6}{6} Theorem \ref{thm:c3-5-1} states, in particular, that the unrestricted limits $\lim_{y\in\Gamma,y\to0}F(x+iy)=F(x)$ exist for almost every $x\in E_n$ whenever $\Gamma$ is the first octant and $F\in H^p(T_\Gamma), p>0$. Corollary \ref{cor:c2-3-23} of Chapter II assures us that the same is true for the iterated Poisson integral of a function $f$ in $L^p(E_n),1<p$. We have already observed that this unrestricted convergence may fail when $f\in L^1(E_n)$; however, if we assume in addition that $|f|(\log|f|)^{n-1}$ (this expression being 0 when $|f|=0$) is locally integrable (which is the case when $f\in L^p(E_n),1<p$) then the iterated Poisson integral of $f$ does converge unrestrictedly to $f(x)$ for almost every $x\in E_n$ (see Jessen, Marcinkiewicz, and Zygmund [1]).

\frp{6}{7} A regular cone $\Gamma \subset E_n$ is called self-dual if its closure $\overline{\Gamma}$ coincides with its dual $\Gamma^*$. An automorphism of a regular cone $\Gamma$ is a linear operator on $E_n$ that maps $\Gamma$ onto $\Gamma$. The collection of all automorphisms of $\Gamma$ is a closed subgroup $\Sigma = \Sigma(\Gamma)$ of the general linear group $GL(E_n)$. If $\Sigma$ is transitive (that is, whenever $x$ and $y$ belong to $\Gamma$ there exists an automorphism $\rho \in \Sigma$ such that $\rho x = y$) and $\Gamma$ is self-dual then it is called a (homogeneous) domain of positivity\index{domain of positivity@domain of positivity}. Except for an exceptional lower-dimensional cone all domains of positivity can be realized as direct sums of the following four types of domains of positivity, often referred to as the ``classical domains'' (for details see Koecher [1], Rothaus [1], and Vinberg [1]):\index{homogeneous domain of positivity@homogeneous domain of positivity}\index{self-dual cone@self-dual cone}

(i) The forward light cone\index{forward light cone@forward light cone}. This is the cone $C_n \subset E_n, n > 1$, consisting of all $x = (x_1, x_2, \ldots, x_n) \in E_n$ such that $x_1^2 - x_2^2 - \cdots - x_n^2 > 0$ and $x_1 > 0$. The automorphism group $\Sigma = \Sigma(C_n)$ consists of all those linear transformations $\rho$ preserving (up to a positive multiplicative constant) the bilinear form $(x, y) = x_1 y_1 - x_2 y_2 - \cdots - x_n y_n$ and the inequality $x_1 > 0$.

(ii) If, for some positive integer $m, n = m(m + 1)/2$ we can identify $E_n$ with the vector space of all $m \times m$ real symmetric matrices. $S_m \subset E_n$ is

then the cone of all positive definite real symmetric matrices. The automorphism group $\Sigma = \Sigma(S_m)$ can be obtained from the real general linear group $GL(E_m)$: to each nonsingular linear transformation $g$ on $E_m$ there corresponds the automorphism $\rho = \rho_g$ of $S_m$ that maps $x \in E_n$ onto $\rho x = g x g^*$ (where $g^*$ is the transpose of $g$). The tube $T_{S_m}$ over this cone is known as the Siegel generalized upper half-plane\index{Siegel upper half-plane@Siegel upper half-plane} (when $m = 1 = n$ this tube is the upper half-plane $E_2^+$).

(iii) When $n = m^2$, $E_n$ can be identified with the real vector space of $m \times m$ complex hermitian matrices (that is, complex $m \times m$ matrices x that are equal to their conjugate transpose $x^*$). $H_m \subset E_n$ is then the cone of all positive definite $m \times m$ complex hermitian matrices. $\Sigma(H_m)$ can be obtained from the complex general linear group $GL(C_m)$: To each nonsingular linear transformation $g$ on $C_m$ there corresponds the automorphism $\rho = \rho_g$ of $H_m$ that maps $x \in E_n$ onto $gxg^*$ where $g^*$ is the conjugate transpose of $g$.

(iv) When $n = 2m^2 - m$, $E_n$ can be identified with the real vector space of $m \times m$ quaternionic hermitian matrices. $Q_m \subset E_n$ is then the cone of all positive definite $m \times m$ quaternionic hermitian matrices. $\Sigma(Q_m)$ can be obtained from the quaternionic general linear group by letting $\rho_g x = g x g^*$ for $x \in E_n$, where $g^*$ is the quaternionic conjugate transpose of the nonsingular $m \times m$ quaternionic matrix $g$.

The Cauchy and Poisson kernels associated with these cones can be calculated explicitly. When $x + iy$ is a point of the tube $T_{C_n}$ whose base is the forward light cone then the Poisson kernel evaluated at this point has the value

\[
\mathcal{P}(x,y)=c_{n}\frac{(y,y)^{n/2}}{\{[(x,x)-(y,y)]^{2}+4(x,y)^{2}\}^{n/2}}=c_{n}\frac{(y,y)^{n/2}}{|(x+iy,x+iy)|^{n}},
\]

when $x + iy \in T_{S_m}$, $\mathcal{P}(x, y) = a_n\{\det y/|\det(x + iy)|^{2\binom{n+1}{2}/2}$, for $x + iy \in T_{H_m}$, $\mathcal{P}(x, y) = a_n'\{\det y/|\det(x + iy)|^{2}\}^n$. To each quaternion $a_0 + ia_1 + ja_2 + ka_3 = (a_0 + ja_2) + i(a_1 + ja_3)$ there corresponds a $2 \times 2$ matrix over the complex numbers

\[
\begin{pmatrix}\alpha&-\bar{\beta}\\ \beta&\bar{a}\end{pmatrix},
\]

where $\alpha = (a_0 + j a_2)$ and $\beta = (a_1 + j a_3)$. This correspondence is an isomorphism and can be extended between $m \times m$ quaternionic hermitian matrices $x$ and $(2m) \times (2m)$ matrices $x'$ with complex coefficients. In terms of this correspondence we have, for $x + iy \in T_{Q_m}$, $\mathcal{P}(x, y) = a_n''(\det y') / |\det(x') + iy')|^{2j(2n-1)/2}$.

\frp{6}{8} The result described in 6.6 when $\Gamma$ is the first octant is not true for general cones. In fact, unrestricted convergence of Poisson integrals may fail even for functions in $L^p(E_n), p > 1$. For example, when $\Gamma = C_n$ (the

forward light cone), $n \geq 3$, and $1 \leq p < \infty$ then there exists an $f \in L^{p}(E_{n})$ such that for almost every $x \in E_{n}$

\[
\lim\sup u(x,y)=\lim\sup\int\limits_{E_{n}}\mathcal{P}(x-t,y)f(t)dt=\infty
\]

as $y \in \Gamma$ tends to 0 unrestrictedly (see Stein and N. J. Weiss [1]). We do obtain a positive result, however, when we limit ourselves to restricted convergence and when the base $\Gamma$ of the tube is one of the domains described in 6.7. In fact, if $f \in L^p(E_n)$, $1 \leq p$, then, for almost every $x \in E_n$, the Poisson integral $u(x, y)$ converges to $f(x)$ as $y \in \Gamma$ tends to 0 restrictedly see Stein and N. J. Weiss [1]).

\frp{6}{9} The classical \textit{Cayley\index{Cayley transform@Cayley transform} transform} is the linear fractional transformation $w = (z - i)/(z + i)$ mapping the upper half-plane onto the unit disk $|w| < 1$. The Siegel upper half-plane corresponding to $S_m$ described in 6.7 (ii) can be mapped by a generalized Cayley transform onto the ``disc'' $D$ of $m \times m$ complex symmetric matrices $w$ satisfying $w^*w < I$, where $I$ is the $m \times m$ identity matrix. This mapping is given by $w = (z - iI) \times (z + iI)^{-1}$ when $z = x + iy \in T_{S_m}$. The ``distinguished boundary'' or ``spine'' $\{z = x + iy \in C_n; y = 0\}$ then corresponds to the symmetric unitary matrices. There are similar transformations of the tubes over the other classical domains introduced in 6.7 (see Bochner [1], Pyatetskii-Shapiro [1] and Korányi and Wolf [1]). The images of these tubes form an important subclass of the bounded symmetric domains of Cartan (see Hua [1] and Helgason [1]).

\frp{6}{10} Let $D$ be the domain of $m \times m$ complex matrices $w$ satisfying $w^*w < I$. We can consider $D$ as a domain in $C_m^2$. We say that a complex valued function $F$ belongs to $H^p(D)$, $p > 0$, if it is holomorphic in $D$ and $\int_U |F(\rho u)|^p du \leq M < \infty$ for $0 \leq \rho < 1$, where $u$ ranges over the unitary group $U$ and $du$ is the element of Haar measure on this group. It can be shown that for $F \in H^p(D)$, $p > 0$, (i) $\lim_{\rho \to 1} F(\rho u) = F(u)$ exists for almost every $u \in U$ and (ii) $\int_U |F(\rho u) - F(u)|^p du \to 0$ as $\rho \to 1$. A similar result holds for any Reinhardt circular region and, thus, for the bounded symmetric domains of Cartan (see Bochner [5]). When $D$ is the polydisc $\{z = (z_1, z_2, \ldots, z_n) \in C_n; |z_j| < 1, j = 1, 2, \ldots, n\}$ the space $H^p(D)$. $p > 0$, consists of those holomorphic functions $F$ defined in $D$ that satisfy $\int_0^{2\pi} \cdots \int_0^{2\pi} |F(r_1 e^{t\theta_1}, \ldots, r_n e^{t\theta_n})|^p \, d\theta_1 \cdots d\theta_n \leq M < \infty$ for $0 \leq r_j < 1$, $j = 1, 2, \ldots, n$. Similarly, when $D$ is the open unit sphere $\{z = (z_1, \ldots, z_n) \in C_n; |z_1|^2 + \cdots + |z_n|^2 < 1\}$ a holomorphic function $F$ defined in $D$ is said to belong to $H^p(D)$ if $\int_{\partial D} |F(\rho z')|^2 \, dz' \leq M < \infty$ for $0 \leq \rho < 1$, where the boundary $\partial D$ of $D$ may be identified with the surface $\Sigma_{2n-1}$ of

the unit sphere in $E_{2n}$ and $dz'$ is then the element of surface area. In these cases we can strengthen the two results (i) and (ii) to include nontangential convergence to the boundary. When D is the polydisc the approach can be ``unrestricted''; that is, the radii $r_1, r_2, \ldots, r_n$ can approach 1 independently (for details see Zygmund [1] and [3]). See also Korányi [1].

\frp{6}{11} Suppose $D$ is the domain $\{z = (z_1, z_2, \ldots, z_n); \mathcal{I}z_1 > \sum_{j=2}^{n} |z_j|^2\}$. The mapping $z \to w = (w_1, w_2, \ldots, w_n)$ given by $w_1 = (z_1 - i)/(z_1 + i)$ and $w_j = 2z_j/(z_j + i)$, $j = 2, \ldots, n$, transforms $D$ into the open unit sphere $\{w \in C_n; \sum_{j=1}^{n} |w_j|^2 < 1\}$ (when $n = 1$ this mapping is the Cayley transform). A general class of domains that contains $D$ as well as the tube domains over convex cones has been introduced by Pyatetskii--Shapiro [1]; these domains are called by him Siegel domains\index{Siegel domains of the second kind@Siegel domains of the second kind} of the second kind (see also Korányi [2]). An $H^p$ theory for such domains generalizing Theorem \ref{thm:c3-5-1} can be found in Stein [6].

\frp{6}{12} Tubes over the $n$-dimensional cube $Q = Q_n = \{y \in E_n; 0 < y_j < 1, j = 1, 2, \ldots, n\}$ share many properties with tubes over the first octant. For example, a result analogous to Theorem \ref{thm:c3-5-1} holds for the former domains: If $F$ belongs to $H^p(T_Q), p > 0$, and $y_0$ is a vertex of $Q$, then $F$ has a nontangential limit $F(x + iy_0)$ in each variable at almost every $x$ in $E_n$. In particular $\lim_{y \in Q_n, y \to y_0} F(x + iy) = F(x + iy_0)$ exists for almost every $x \in E_n$. Moreover,

\[
\int\limits_{E_{n}}|F(x+iy)-F(x+iy_{0})|^{p} dx
\]

tends to zero as $y \in Q$ tends to $y_0$. The proof is similar to that of Theorem \ref{thm:c3-5-1}. The rôle of the upper half-plane is assumed by the band $\{z = x + iy \in C_1; x \in E_1, 0 < y < 1\}$ and the Poisson integral of a function $f$ defined on the boundary of this region is given by the formula

\[
u(x+i y)=\textstyle\frac{1}{2}\sin\pi y\int\limits_{-\infty}^{\infty}\left[\frac{f(x-t)}{\cosh\pi t-\cos\pi y}+\frac{f(x-t+i)}{\cosh\pi t+\cos\pi y}\right]d t
\]

(see comments following the proof of Lemma \ref{lem:c5-4-2} in Chapter V). With the aid of an appropriate linear transformation we can extend this result to $H^p$ spaces of functions defined on tubes with polyhedral bases. We then obtain the following extension of Theorem \ref{thm:c3-5-5} (and Theorem \ref{thm:c3-5-5p}):

If $F$ belongs to $H^p(T_B)$, $p > 0$, where $B$ is an open convex subset of $E_n$, and $y_0$ is a point on the boundary of $B$ then the restricted limits

\[
\lim_{y\to y_{0},y\in B}F(x+iy)=F(x+iy_{0})
\]

exist for almost every $x \in E_n$ and $\int_{E_n} |F(x + iy) - F(x + iy_0)|^p dx \to 0$ as $y \in B$ tends restrictedly to $y_0$.

\frp{6}{13} In the third section of this chapter we introduced a Poisson kernel associated with each regular cone $\Gamma \subset E_n$. This was done by extending to $n$-dimensions a relationship between the classical Poisson and Cauchy kernels associated with the upper half-plane (i.e., the tube whose base $\Gamma \subset E_1$ consists of the positive reals). Other relations of this kind exist. Perhaps the simplest is expressed by the following equality:

\[
\begin{aligned}K(x+iy)&=-\frac{1}{2\pi i(x+iy)}=\frac{1}{2}\left\{\frac{y}{\pi(x^{2}+y^{2})}+i\frac{x}{\pi(x^{2}+y^{2})}\right\}\\&=\frac{1}{2}\left\{P(x,y)+i Q(x,y)\right\}.\end{aligned}
\]

Consequently, $P(x,y)=2\text{Re}\{K(x+iy)\}$. The kernel $Q(x,y)=2\mathcal{I}\{K(x+iy)\}$ is called the conjugate Poisson\index{conjugate Poisson kernel@conjugate Poisson kernel!for cones@for cones} kernel. Motivated by this relationship, therefore, one might be led to consider

\[
P(x,y)=2\text{Re}\left\{K(x+iy)\right\}=2\text{Re}\left\{\int\limits_{\Gamma^{*}}e^{2\pi i(x+iy)\cdot t}dt\right\}
\]

and

\[
Q(x,y)=2\mathcal{I}\{K(x+i y)\}=2\mathcal{I}\Biggl\{\intop_{\Gamma^{*}}e^{2\pi i(x+i y)\cdot t}d t\Biggr\}
\]

as generalizations of the Poisson and conjugate Poisson kernels associated with a regular cone $\Gamma \subset E_n$. For example, when $\Gamma$ is the first quadrant of $E_2$,

\[
P(x_{1},x_{2};y_{1},y_{2})=P(x,y)=\frac{1}{\pi^{2}}\frac{y_{1}y_{2}-x_{1}x_{2}}{(x_{1}^{2}+y_{1}^{2})(x_{2}^{2}+y_{2}^{2})}.
\]

In particular, we see that $P(x, y)$ is not necessarily a positive kernel. For this and as well as other reasons this kernel is not a satisfactory approximation to the identity. Nevertheless, it can be used together with $Q(x, y)$ in order to see how the real and imaginary parts of an $H^{2}$ function are related to each other and to their boundary values.

If $\Gamma \subset E_n$ is a regular convex open cone we let $\chi^*$ denote the characteristic function of the dual cone $\Gamma^*$ and sgn* the signum\index{signum function@signum function!for cones@for cones} function associated with $\Gamma^*$: sgn* $t = 1$ if $t \in \Gamma^*$, sgn* $t = -1$ if $-t \in \Gamma^*$ and sgn* $t = 0$ otherwise. It then follows immediately from the definitions above that

\[
P(x,y)=\int\limits_{E_{n}}e^{2\pi i x\cdot t}e^{-2\pi|y\cdot t|}[\chi^{*}(t)+\chi^{*}(-t)]d t
\]

and

\[
\begin{aligned}{Q(x,y)}&{{}=-i\intop_{E_{n}}e^{2\pi i x\cdot t}e^{-2\pi|y\cdot t|}[\chi^{*}(t)-\chi^{*}(-t)]\;d t}\\ {}&{{}=-i\intop_{E_{n}}e^{2\pi i x\cdot t}e^{-2\pi|y\cdot t|}\text{sgn*}\,t\;d t.}\\ \end{aligned}
\]

Thus, for each $y \in \Gamma$, $\hat{P}_y(t) = [\chi^*(t) + \chi^{*}(-t)] e^{-2\pi|y \cdot t|} \text{ and } \hat{Q}_y(t) = (-i \text{sgn*}\,t) e^{-2\pi|y \cdot t|} \text{ are the values of the Fourier transforms of } P_y = P(\cdot, y) \text{ and } Q_y = Q(\cdot, y)$.

Now suppose $F = u + iv \in H^2(T)$. By Corollary \ref{cor:c3-3-4} the $L^2$ boundary values $\lim_{y \in \Gamma, y \to 0} F(x + iy) = u(x) + iv(x)$ exist. It is not hard to establish the following results:

(i) For all $z \in T_{\Gamma}$, $F(z) = \int_{E_n} 2K(z - \xi) u(\xi) \, d\xi$;

(ii) $v(z) = v(x + iy) = \int_{E_n} Q(x - \xi, y)u(\xi) \, d\xi$, when $z \in T_\Gamma$, and the boundary values $v(x)$ and $u(x)$ satisfy the relationship $\hat{v}(t) = (-i \text{sgn*}\,t)\hat{u}(t)$;

(iii) A real-valued function $u \in L^2(E_n)$ is the real part of the $L^2$ boundary value $F(x) = \lim_{y \in \Gamma, y \to 0} F(x + iy)$ of an $F \in H^2(T_\Gamma)$ if and only if $u(t)$ vanishes for almost every $t$ outside $\Gamma^* \cup (-\Gamma^*)$, where $-\Gamma^* = \{t \in E_n; -t \in \Gamma^*\}$.

The first part of (ii) follows immediately from (i) by equating imaginary parts. The second part is then obtained by taking the Fourier transform of both sides and passing to the limit as $y \in \Gamma$ tends to 0. If $u \in L^2(E_n)$ is the real part of the boundary value of an $F \in H^2(T_r)$ then from (i) and the definition of the Cauchy kernel we obtain

\[
F(z)=2\int\limits_{E_{n}}K(z-\xi)u(\xi)\;d\xi=2\int\limits_{\Gamma^{*}}\hat{u}(t)e^{2\pi i z\cdot t}\;d t
\]

for all $z \in T_{\Gamma}$. We note that, in any case, an $F$ defined in this way belongs to $H^2(T_{\Gamma})$ (this follows from Theorem \ref{thm:c3-3-1}) and its real part $u_y(x) = u(x + iy)$ equals $\int_{E_n} P(x - \xi, y)u(\xi) \, d\xi$. Thus, taking Fourier transforms

\[
\hat{u}_{y}(t)=e^{-2\pi|y\cdot t|}(\chi^{*}(t)+\chi^{*}(-t))\hat{u}(t).
\]

Part (iii) is an immediate consequence of this equality. Part (i) is obtained by adding the left and right sides of the equalities

\[
F(z)=\int\limits_{E_{n}}K(z-\xi)F(\xi)d\xi
\]

and $0 = \int_{E_n} K(z - \xi) \overline{F(\xi)} \, d\xi$. The first of these identities follows from

Theorem \ref{thm:c3-3-6}. The second is an easy consequence of Theorem \ref{thm:c3-3-1} and the definition of the Cauchy kernel.

\frp{6}{14} In the case when $n=1$ and $\Gamma$ is the positive real axis then every real-valued $u\in L^{2}(E_{1})=L^{2}(-\infty,\infty)$ is the real part of the boundary value of an $F=u+iv\in H^{2}(E_{2}^{+})$. This is an immediate consequence of 6.13, (iii) and the fact that $\Gamma^{*}\cup(-\Gamma^{*})=(-\infty,\infty)=E_{1}$. From (ii) and Plancherel's theorem we see that the mapping $u\to v$ is well defined and is an isometry of $L^{2}(E_{1})$ (in fact, it is a unitary transformation). This mapping is called the Hilbert transform\index{Hilbert transform@Hilbert transform}. We shall show later in Chapter V, Section 2, that the Hilbert transform can be defined on $L^{p}(E_{1})$ and is a bounded linear transformation into $L^{p}(E_{1})$, $1<p<\infty$. This fact is equivalent to saying that there exists a constant $A_{p}<\infty$ such that, whenever $u+iv$ is the boundary value of an $F\in H^{p}(E_{2}^{+})$,

\[
\|v\|_{p}=\bigg(\intop_{-\infty}^{\infty}|v(x)|^{p}d x\bigg)^{1/p}\leq A_{p}\bigg(\intop_{-\infty}^{\infty}|u(x)|^{p}d x\bigg)^{1/p}=A_{p}\|u\|_{p}.
\]

This inequality is known as the M. Riesz inequality. This one-dimensional result can be extended to $n$ dimensions:

If $F = u + iv \in H^p(T_\Gamma)$, $1 < p < \infty$, where $\Gamma$ is an open convex cone in $E_n$, then there exists a constant $A_p < \infty$ such that for $y \in \Gamma$, and for $y = 0$ (this limiting value can be taken in the sense of Theorem \ref{thm:c3-5-6})

\[
\biggl(\int_{E_{n}}|v(x+i y)|^{p}d x\biggr)^{1/p}\leq A_{p}\biggl(\int_{E_{n}}|u(x+i y)|^{p}d x\biggr)^{1/p}.
\]

By a change of coordinates we can reduce this to the case where the ray $(\eta_1, 0, \ldots, 0)$, $\eta_1 > 0$, belongs to $\Gamma$. If $y \in \Gamma$ is fixed and $\eta = (\eta_1, 0, \ldots, 0)$, $\eta_1 > 0$, then $y + \eta \in \Gamma$ and, by the one-dimensional result,

\[
\begin{aligned}&\int\limits_{-\infty}^{\infty}|v(x_{1}+iy_{1}+i\eta_{1},\ldots,x_{n}+iy_{n})|^{p} dx_{1}\\&\leq A_{p}^{p}\int\limits_{-\infty}^{\infty}|u(x_{1}+iy_{1}+i\eta_{1},\ldots,x_{n}+iy_{n})|^{p} dx_{1}.\end{aligned}
\]

Integrating both sides with respect to $x_2, \ldots, x_n$ and then letting $\eta_1 \to 0$ we obtain the desired inequality. The case $y = 0$ follows from this one by letting $y \in \Gamma$ tend to 0.

\frp{6}{15} Let $K$ be any open convex set containing no (entire straight) lines. Suppose $p \notin K$ and let $\Gamma$ be the cone generated by $p$ and $K$. Then $\Gamma$ does not contain any (entire straight) lines. To see this assume first that $n = 2$.

Consider the supporting lines to $K$. Either they are all parallel (which would imply that $K$ contains such a line) or there are two properly intersecting ones. Their intersection determines a cone $\Gamma_0$ containing $K$ such that no such line lies within $\Gamma_0$. The required property for $\Gamma$ is now obvious. Assume now that $n > 2$, and that $l$ is a line lying in $\Gamma$. Let $\pi$ be a plane containing $l$ and $p$ and let $K' = K \cap \pi$, $\Gamma' = \Gamma \cap \pi$. Then $K'$ is open, convex and does not contain any line; moreover, $\Gamma'$ is the cone generated by $K'$ and $p$, and contains the line $l$, which contradicts the case $n = 2$, already proved.

\frp{6}{16} If the convex open cone $\Gamma$ contains no straight lines, then $\Gamma$ is regular. It follows from Theorem \ref{thm:c3-3-1} that $H^2(T_r)$ contains a function which is not identically zero. To prove the assertion about $\Gamma$ one uses the fact (see Rockafellar [1], Theorem 14.1) that $\bar{\Gamma} = (\Gamma^*)^*$. If $\Gamma^*$ had no interior it would be contained in a hyperplane, which would imply that $\bar{\Gamma}$ contained a straight line, and hence $\Gamma$ would contain a straight line.

\bibnoteshead

The basic results of the classical theory of $H^{p}$ spaces can be found in Zygmund [1], Chapter VII. The $H^{2}$ theory over tubes was first studied by Bochner to whom the basic result, Theorem \ref{thm:c3-2-3}, is due (see Bochner [8]). Many of the results included in Sections 2, 3 and 5 are presented in Stein, G. Weiss and M. Weiss [1]. The Cauchy kernel associated with a convex cone was introduced by Bochner (see [1]), he also calculated this kernel explicitly for cones corresponding to the classical domains. For other extensions of the Paley--Wiener theorem to n-dimensions see Plancherel and Polya [1] and Stein [1]. The notion of curvilinear restricted convergence described in \eqref{frp:c3-6-5} was developed by Calderón and Zygmund [1]. The basic results concerning the Hilbert transform are presented in Titchmarsh [2] and Zygmund [1]. For earlier papers dealing with $H^{p}$ theory of several complex variables see Bochner [2], Zygmund [3] and Bochner [5]. For an introduction to the theory of convex sets see Valentine [1], and Rockafellar [1].

